%% Pre-registration of the TESS Sector 110 test of TIC 4206066: main part (Sections 1 to 7) and the full protocol
%% (Appendices A to N). Same class and conventions as the author's preprint (aastex702, default style).
%% Rules, numbers and code are those of the frozen deposit tic4206066_S110_prereg_final (SHA256SUMS sha256 310cb93d...);
%% RULE_MAP.md maps every rule of the earlier Markdown text to its place here.
\documentclass{aastex702}
\usepackage{enumitem}
\usepackage{longtable}
\usepackage{booktabs}
\usepackage{array}

\defcitealias{GaiaMission}{Gaia Collaboration et al. 2016}
\defcitealias{GaiaDR3}{Gaia Collaboration et al. 2023}
\defcitealias{DESY3}{DES Collaboration et al. 2022}
\defcitealias{Astropy2013}{Astropy Collaboration et al. 2013}
\defcitealias{Astropy2018}{Astropy Collaboration et al. 2018}
\defcitealias{Astropy2022}{Astropy Collaboration et al. 2022}
\defcitealias{Lightkurve2018}{Lightkurve Collaboration et al. 2018}
\defcitealias{COSPrereg}{Center for Open Science 2026}
\shorttitle{Pre-registered Sector 110 Follow-up of TIC 4206066}
\shortauthors{Rabtsevich}
\journalinfo{Pre-registration, written before the TESS Sector 110 observations}

\newcommand{\ddet}{d_{\rm det}}
\newcommand{\dlo}{d_{\rm lo}}
\newcommand{\dhi}{d_{\rm hi}}
\newcommand{\Tfourteen}{T_{14}}
\newcommand{\code}[1]{\texttt{#1}}
\let\origunderscore\_
\renewcommand{\_}{\origunderscore\allowbreak}
\newcommand{\slashbreak}{/\allowbreak}
\newcommand{\Sens}{S}
\newcommand{\runin}[1]{\par\smallskip\noindent\textbf{#1}\enspace}
\setlist{leftmargin=*,itemsep=1pt,topsep=3pt}
\setlist[itemize,2]{label=$\circ$}

\begin{document}

\title{A Pre-registered TESS Sector 110 Follow-up Protocol for Two Transit-like Signals in TIC 4206066}

\author[gname=Pavel,sname=Rabtsevich]{Pavel Rabtsevich}
\affiliation{Independent Researcher, Tarragona, Spain}
\email{pavel.rabtsevich.research@gmail.com}

\begin{abstract}
Archival photometry from the Transiting Exoplanet Survey Satellite (TESS) of the nearby K dwarf TIC 4206066 shows two
transit-like signals: signal A with a period of 3.18\,d and a tentative signal B with a period of 11.13\,d and less archival support.
Both were selected after the data had been inspected, so their archival significance cannot be taken at face value.
TESS observes the star again in Sector 110, from 2026 October 31 to November 26, at 2-minute cadence. These data did not
exist when this protocol was written, so they give an independent test. The protocol fixes the predicted transit
times, the expected depths, the processing of the light curve, a rule that gives each signal one of four outcomes
(recovered, inconsistent, excluded, or insufficient data), a timing rule, and a pixel-level check that signal A comes
from TIC 4206066. Frozen code computes the outcomes from valid inputs; this text defines which inputs and runs are
valid. A recovery would not show that a signal is caused by a planet, and an undecided result would not show that it
is not.
\end{abstract}

\keywords{TESS, TIC 4206066, exoplanet candidates, transit photometry, pre-registration}

\section{Question and status of the data} \label{sec:question}

\runin{Question.} Do the two transit-like signals seen in archival TESS \citep{Ricker2015} photometry of TIC 4206066
reappear in new data, at the times and depths predicted from the archival data? For signal A, the protocol also asks
whether the dip again comes from the pixels of TIC 4206066 rather than from a neighbouring star.

\runin{Status of the data.} The archival TESS Sectors 6, 32 and 98 were inspected before this protocol was written:
the signals were found in them, and the predictions are fitted to them. This is therefore not a blind pre-registration
of the whole study. It pre-registers one future test. TESS observes TIC 4206066 again in Sector 110, from 2026 October
31 to November 26 (camera 1, CCD 2), at 2-minute cadence, through an approved Director's Discretionary Target request
(DDT 100, listed at \url{https://tess.mit.edu/science/ddt/}). No Sector 110 data existed when this protocol was
written, and they play no part in finding or fitting the signals. They therefore test the signals with independent
noise, though on the same star field, provided that every choice of the analysis is fixed before the data arrive.
This protocol counts as a pre-registration only if it is deposited before 2026 October 31.

\runin{What kind of document this is.} Publishing an ephemeris that predicts future transits is routine in exoplanet
work; for example, \citet{IvshinaWinn2022} published updated ephemerides of 382 planets for scheduling follow-up
observations. A pre-registration adds that the analysis and the decision rules are fixed in advance as well. This
document is a pre-registration in the sense of the Center for Open Science: a time-stamped plan, made public before the
data exist, that keeps the planned (confirmatory) analysis apart from any later, exploratory one
\citepalias{COSPrereg}. Its main part (Sections~\ref{sec:question} to \ref{sec:protocol}) follows the questions of
short pre-registration forms, for example the AsPredicted entry of \citet{AsPredicted13289}: which data were already
seen, what is measured, how it is analysed, what is excluded, and what decides the result. The full rules are in
Appendices~\ref{app:data} to \ref{app:codemap}, and Section~\ref{sec:protocol} states which text governs. It is not a
blind analysis, in which data that already exist are analysed with the result hidden until the choices are fixed
(e.g., \citetalias{DESY3}, their Appendix D; \citealt{Asgari2021}, their Appendix F), and it is not a Registered
Report, whose protocol is peer reviewed before the data are taken.

\section{How the two signals were found, and why they are tested} \label{sec:signals}

\runin{The search.} I searched the Sector 98 light curves of the MIT Quick-Look Pipeline \citep[QLP;][]{Huang2020} of
all 126,246 stars with TESS magnitude $T \leq 11.5$ for periodic transit-like dips, using box least squares after
detrending. A signal had to reach a multiple-event statistic (MES; a signal-to-noise ratio of the folded transit that
allows for correlated noise, as computed by the automated vetting tool LEO-vetter; \citealt{Kunimoto2025}) of at least
9. The search returned 23,949 signals on 16,300 stars. After automated vetting, 141 signals on 140 stars remained; 70
of them with MES $\geq 15$ formed the main list, and 71 with MES $< 15$ formed an exploratory list. TIC 4206066 (Gaia
DR3 3220388198192519424; a K dwarf at 35.6\,pc) is the only star with two signals among the 141: A, with a period of
3.18\,d and MES 10.6, and B, with a period of 11.13\,d and MES 11.0. Both are in the exploratory list, and both passed
all automated false-alarm and false-positive tests of the vetting.

\runin{Why the archival evidence is not enough.} The exploratory list (MES $< 15$) was added to the search protocol on
2026 September 25, about two hours after signal B of this star had been found, so it was defined after results for
individual stars, including this one, had been seen. I did not record why this star was the first to be examined in
detail. No false-alarm rate or completeness was computed for the search as a whole. The archival significance of A
and B therefore cannot be taken at face value. A test on data that did not exist when the signals were chosen avoids
this problem, and that is the purpose of this protocol.

\runin{What the archival data show.} Signal A recurs in the older TESS Sectors 6 and 32. Its expected box depth from
those two sectors is $442 \pm 52$\,ppm (parts per million of the stellar flux; Table~\ref{tab:depths}). The transit fit of the preprint (radius ratio
$0.0204 \pm 0.0010$ on a star of $0.623 \pm 0.019\,R_\odot$) corresponds to a planet of $1.4 \pm 0.1$ Earth radii if
it orbits TIC 4206066. Archival difference images place the dip of A about $5''$ east of the target, at $(+5 \pm 3, 0
\pm 3)''$, an offset of about $2\sigma$ from TIC 4206066 that does not exclude it. A dry run of the frozen localisation code on the
Sector 32 pixels (a run on archival data, which never produces a label) gives an empirical $p$ of 0.156, which would
read ``compatible'' under the rule of Section~\ref{sec:localisation} (Appendix~\ref{app:validation}).

Signal B is tentative. In Sector 98 alone it is as strong as A (MES 11.0 against 10.6), but its archival support is
weaker. Under the counting rule of this protocol it has two transits in Sector 98 and one in Sector 6; its two Sector
32 events fall next to transits of A and cannot be counted (Table~\ref{tab:depths}). Its expected depth, $586 \pm
137$\,ppm, therefore rests on the single Sector 6 transit. The archival dips of B were found only at the times
predicted from the Sector 98 ephemeris, with a test defined after the data had been inspected; an earlier search that
was not anchored on that ephemeris did not find them. A weaker feature with a period of 1.76\,d (about 130\,ppm) is
also seen in Sector 98; it is not tested here (Section~\ref{sec:limits}).

The discovery and the checks made so far are described in a preprint (version DOI \href{https://doi.org/10.5281/zenodo.22967456}{doi:10.5281/zenodo.22967456}) with
its data and code (concept DOI \href{https://doi.org/10.5281/zenodo.22967098}{doi:10.5281/zenodo.22967098}, all versions). The preprint claims no statistical
validation; the main alternative it leaves open for signal A is a planet that transits an unresolved companion star
bound to TIC 4206066, rather than TIC 4206066 itself.

\section{Predictions} \label{sec:predictions}

\subsection{Ephemerides} \label{sec:ephem}

All times used by the analysis are Barycentric TESS Julian Dates, BTJD $=$ BJD$_{\rm TDB} - 2457000$. The ephemerides
are in the file \code{ephemeris\_final.json}:
\begin{itemize}
\item A: period $P = 3.1827844492913377$\,d, reference mid-time $T_0 = 3998.7348271847636$ BTJD, transit duration (first to
  last contact) $\Tfourteen = 1.8705$\,h, with the covariance of $(T_0, P)$ in the file ($\sigma_{T_0} = 0.00243$\,d $= 3.5$\,min,
  $\sigma_P = 0.0000051$\,d, correlation 0.70);
\item B: $P = 11.132745380039134$\,d, $T_0 = 4000.8658739613074$ BTJD, $\Tfourteen = 1.6581$\,h, unchanged since the
  first internal draft of this protocol; its prediction uncertainties (Table~\ref{tab:timesB}) are 1.5 times the
  per-epoch uncertainties of the joint fit, because its timing errors are not calibrated by injections.
\end{itemize}
The values are given with the digits stored in the file. The ephemeris of A rests on one choice made after looking at
the archival data: in two Sector 32 transits of A, signal B falls close to A, and B is modelled jointly with A instead
of being masked. Masking B instead would shift the predicted times of A by 0.4\,min (2.6 instead of 3.0\,min earlier
than in the internal draft v2, Appendix~\ref{app:history}), well inside the timing tolerance of at least 12.3\,min
(Section~\ref{sec:timing}).

\subsection{Predicted transits in Sector 110} \label{sec:times}

Tables~\ref{tab:timesA} and \ref{tab:timesB} list every Sector 110 epoch, with the prediction uncertainty $\sigma$ of
each. The scripts use BTJD only. Two calendar columns are added for convenience. ``BJD$_{\rm TDB}$ calendar'' is the
same barycentric time written as a UTC calendar date (the \code{tc\_utc\_calendar} strings of
\code{ephemeris\_final.json}); it is \emph{not} the time at which the transit is seen. ``UTC at Earth'' is the arrival
time at the centre of the Earth in UTC, computed with astropy from the Gaia DR3 position \citepalias{GaiaMission,GaiaDR3}
with the light-travel time to the solar-system barycentre removed; at TESS the arrival time is within 2\,s of it. In
Sector 110 the arrival time is 5.6 to 7.2\,min earlier than the BJD$_{\rm TDB}$ calendar date, because the star is
then on the Earth's side of the Sun.

\begin{deluxetable}{rcccc}
\tablecaption{Predicted mid-transit times of signal A in Sector 110.\label{tab:timesA}}
\tablehead{\colhead{Epoch} & \colhead{BTJD} & \colhead{BJD$_{\rm TDB}$ calendar} & \colhead{UTC at Earth} &
  \colhead{$\sigma$ (min)}}
\startdata
109 & 4345.65833 & 2026-11-01 03:46:51 & 2026-11-01 03:41:16 & 4.10 \\
110 & 4348.84112 & 2026-11-04 08:10:03 & 2026-11-04 08:04:13 & 4.11 \\
111 & 4352.02390 & 2026-11-07 12:33:16 & 2026-11-07 12:27:10 & 4.11 \\
112 & 4355.20669 & 2026-11-10 16:56:28 & 2026-11-10 16:50:09 & 4.12 \\
113 & 4358.38947 & 2026-11-13 21:19:41 & 2026-11-13 21:13:10 & 4.12 \\
114 & 4361.57225 & 2026-11-17 01:42:54 & 2026-11-17 01:36:11 & 4.13 \\
115 & 4364.75504 & 2026-11-20 06:06:06 & 2026-11-20 05:59:13 & 4.13 \\
116 & 4367.93782 & 2026-11-23 10:29:19 & 2026-11-23 10:22:17 & 4.14 \\
117 & 4371.12061 & 2026-11-26 14:52:31 & 2026-11-26 14:45:23 & 4.15 \\
\enddata
\end{deluxetable}

\begin{deluxetable}{rcccc}
\tablecaption{Predicted mid-transit times of signal B in Sector 110.\label{tab:timesB}}
\tablehead{\colhead{Epoch} & \colhead{BTJD} & \colhead{BJD$_{\rm TDB}$ calendar} & \colhead{UTC at Earth} &
  \colhead{$\sigma$ (min)}}
\startdata
31 & 4345.98098 & 2026-11-01 11:31:28 & 2026-11-01 11:25:51 & 4.06 \\
32 & 4357.11373 & 2026-11-12 14:42:37 & 2026-11-12 14:36:10 & 4.08 \\
33 & 4368.24647 & 2026-11-23 17:53:46 & 2026-11-23 17:46:44 & 4.11 \\
\enddata
\end{deluxetable}

Only these epochs can enter the Sector 110 test. If the light curve covers at least one frozen epoch of a signal, every
frozen epoch of that signal is listed in the output, observed or not. The closest A and B events are 7.4\,h apart
(centre to centre), so their transit windows never overlap. Nine transits of A are predicted, but at most eight can be
counted in practice. A transit is counted only if it lies outside the first and last 9\,h of each continuous stretch of
data, where the detrending is unreliable (the edge zones, Section~\ref{sec:test}). Epoch 117 needs data until about BTJD
4371.53 (2026-11-27 00:41 UTC at Earth) to clear the edge zone at the end of the light curve, which is after the
scheduled end of the sector; epoch 109 needs continuous data from about BTJD 4345.24 (2026-10-31 17:45 UTC at Earth).

\subsection{Expected depths} \label{sec:depths}

The expected depths are box depths: 1 minus the mean normalised flux inside the predicted transit window. A box depth
is lower than a fitted transit depth. They were computed on the archival sectors with exactly the statistic, counting
rule, detrending and noise model of the test (Section~\ref{sec:test}; file \code{calib/expected\_depths.json}). The error
of each value is the standard deviation of the mean of $n$ windows of that sector's noise pool (Section~\ref{sec:stat}). Notes on the individual
values are in Appendix~\ref{app:depthnotes}. The products are the 2-minute light curves of the Science Processing Operations
Center \citep[SPOC;][]{Jenkins2016}, its full-frame-image light curves \citep[TESS-SPOC;][]{Caldwell2020}, and QLP.

\begin{table}[ht]
\caption{Expected box depths from the archival sectors, with the number $n$ of counted transits and the transits that
were not counted. Epochs are counted from $T_0$ (Section~\ref{sec:ephem}); negative epochs precede it. $D$, $D_{98}$ and
$\dhi$ are used as described in Section~\ref{sec:stat}.\label{tab:depths}}
\footnotesize
\setlength{\tabcolsep}{2pt}
\begin{tabular}{>{\raggedright\arraybackslash}p{0.24\textwidth}>{\raggedright\arraybackslash}p{0.40\textwidth}>{\raggedright\arraybackslash}p{0.28\textwidth}}
\toprule
Sector, product & A: depth, ppm & B: depth, ppm \\
\midrule
S6, TESS-SPOC 30-min & $444.6 \pm 64.7$ ($n = 4$; not counted: epoch $-795$, edge zone; epoch $-789$, coverage below
  80\,\%) & $585.5 \pm 137.3$ ($n = 1$) \\
S32, SPOC 2-min & $436.9 \pm 87.0$ ($n = 6$; epochs $-573$ and $-566$ not tested: their masks join the B masks into a
  break of the detrending) & none counted (both B events sit next to A, same reason) \\
S98, QLP 200-s (discovery sector) & $450.0 \pm 47.8$ ($n = 8$; not counted: epoch $-1$, no baseline on one side) &
  $994.5 \pm 98.1$ ($n = 2$) \\
\midrule
Primary expectation $D$ (S6 and S32, inverse-variance mean) & $441.9 \pm 51.9$ & $585.5 \pm 137.3$ (S6 only, one
  transit) \\
$D_{98}$ (S98; the centre used to build $\dhi$, not $\dhi$ itself) & $450.0 \pm 47.8$ & $994.5 \pm 98.1$ \\
\bottomrule
\end{tabular}
\end{table}

Because both signals were found in Sector 98, their depth there may be biased upward by the selection. The Sector 98
depth $D_{98}$ is therefore used only for the upper edge of the consistency band; the primary expectation $D$ comes from
Sectors 6 and 32. A sector without a counted transit drops out of $D$; for B, the primary expectation $D$ therefore
rests on a single Sector 6 transit, and its error reflects this.

\section{The test} \label{sec:test}

\subsection{Data and processing}

The primary is the first released SPOC 2-min light curve, column \code{PDCSAP\_FLUX} (corrected flux), quality mask
17087. Only if it is absent on 2027 February 28 may the primary use QLP 200-s (mask 1610619071), then TESS-SPOC
200-s (\code{PDCSAP\_FLUX}, mask 17087) if QLP is also absent, using the version public on that date and a saved
MAST search record (Appendix~\ref{app:data}). QLP uses \code{SYS\_RM\_FLUX}, or \code{DET\_FLUX} only if the former
column is absent (Appendix~\ref{app:commands}). An existing primary with no counted transit gives O4; it never
triggers a product switch.

The biweight detrending filter has window $W=0.75$\,d \citep{Hippke2019}; both signals are masked over
$\pm1.5\,\Tfourteen$. Upward outliers above $4\times1.4826$ times the median absolute deviation are removed.
Segments split at gaps longer than 0.3\,d in the cleaned series with both transit masks removed. The first and last
$W/2$ (9\,h) of each segment are edge zones. A transit is counted only if (i) at least 80\,\% of its window has
cadences after cleaning, (ii) at least 2 cadences lie on each side, 0.75 to 3\,$\Tfourteen$ from its predicted
mid-time, (iii) its whole window lies within one segment outside the edge zones, and (iv) it is a frozen Sector 110
epoch. Otherwise it is ``not tested'' and never counts against the signal (Appendix~\ref{app:processing}).

\subsection{Statistic and thresholds} \label{sec:stat}

The depth of one transit is 1 minus the mean detrended flux inside $|t - t_c| < \Tfourteen/2$, where $t_c$ is its
predicted mid-time. The test statistic $d$ is the mean of these depths over the $n$ counted transits. Nothing is
searched in period or in time.

The noise is measured on the same light curve. The same window is placed at fake transit times between the real ones,
and each fake ephemeris goes through the full detrending, masking and counting; the windows obtained form the noise
pool (Appendix~\ref{app:noise}). The pool is resampled 20\,000 times for each of 20 seeds; the medians of the resulting thresholds and sensitivity
are used (Appendix~\ref{app:thresholds}).
From these resamples, for the $n$ actually counted:
\begin{itemize}
\item the detection threshold $\ddet$ is the 99.9th percentile of the mean of $n$ noise windows;
\item the band of depths consistent with the prediction runs from $\dlo$, the 0.1st percentile of noise plus a depth
  drawn from $N(D, D_{\rm err})$, to $\dhi$, the 99.9th percentile of noise plus a depth drawn from
  $N(D_{98}, D_{98,\rm err})$;
\item the sensitivity $\Sens = P(\mathrm{detect} \mid \mathrm{real})$ is the fraction of draws of noise plus
  $N(D, D_{\rm err})$ above $\ddet$.
\end{itemize}
The percentile is not the false-alarm level: because neighbouring noise windows overlap, the bootstrap underestimates
the tails. The false-detection levels were therefore measured on archival data with a separate set of fake ephemerides:
0.47\,\% for A and 0.49\,\% for B overall in SPOC 2-min data, with the level for each number of counted transits in
Table~\ref{tab:fa}. These trials are not independent of each other, so their binomial bounds are approximate.

\subsection{Four outcomes} \label{sec:outcomes}

Each signal receives exactly one outcome (Table~\ref{tab:outcomes}).

\begin{longtable}{>{\raggedright\arraybackslash}p{0.17\textwidth}>{\raggedright\arraybackslash}p{0.43\textwidth}>{\raggedright\arraybackslash}p{0.32\textwidth}}
\caption{Outcome of each signal. All quantities are computed for the $n$ counted transits.\label{tab:outcomes}}\\
\toprule
Outcome & Condition & Meaning \\
\midrule
\endfirsthead
\toprule
Outcome & Condition & Meaning \\
\midrule
\endhead
\textbf{O1 recovered} & $d > \ddet$, $\dlo \leq d \leq \dhi$, and the timing is consistent (Section~\ref{sec:timing}) &
  a dip consistent with the prediction is present \\
\textbf{O2 inconsistent} & $d > \ddet$, and at least one of: $d < \dlo$; $d > \dhi$; the timing is inconsistent & a dip is present, but
  not the predicted one (the failed condition is recorded) \\
\textbf{O3 excluded} & $d \leq \ddet$, $d < \dlo$, and $\Sens \geq 0.95$ & the predicted signal is not present at the
  predicted depth, and the data had at least 95\,\% sensitivity to it \\
\textbf{O4 insufficient} & anything else: fewer than 10 usable cadences, or no registered QLP flux column
  (Appendix~\ref{app:commands}); $n = 0$; a noise pool of fewer than 100 windows; or $d \leq \ddet$ and either $d \geq \dlo$ or
  $\Sens < 0.95$ & the data cannot decide \\
\bottomrule
\end{longtable}

Transits that are not tested reduce $n$; they never count against the signal. O1 and O2 do not depend on the
sensitivity: a $d$ above $\ddet$ gives O1 or O2 whatever the sensitivity; its false-detection level is the one in
Table~\ref{tab:fa} for that $n$. The sensitivity matters only when $d \leq \ddet$. In the dry run on the archival
Sector 32 data the sensitivity for A was 0.92, below the 0.95 needed for O3 (Appendix~\ref{app:validation}). Even a
clean non-detection of A may therefore give O4. For B, with two or three transits, the sensitivity may be far below
0.95. In that case an O1 of B is still possible but an O3 is not, and B will most likely end as O4 unless Sector 110 is
quieter than Sector 32.

O4 means only that the data cannot decide; it says nothing against the signal. O1 means that a dip consistent with the prediction is
present in new data; it does not show that the signal comes from a planet (Section~\ref{sec:limits}).

\subsection{Timing rule} \label{sec:timing}

The timing is checked only when $d > \ddet$.
\begin{itemize}
\item The common time shift $\Delta t$ of all counted transits is the shift between $-30$ and $+30$\,min, in steps of
  0.5\,min, that maximises their mean depth. If several shifts share the maximum (a plateau, because the grid is finer
  than the cadence), the middle one is taken; for an even number of them, the later of the two middle shifts.
\item $\sigma_{\rm meas}$ is the standard deviation of the recovered shift over 300 injections of a box of depth $d$ into
  $n$ random noise-pool windows, with a true shift of 0 and a fixed seed.
\item $\sigma_{\rm pred}$ is the mean of the prediction uncertainties (Tables~\ref{tab:timesA} and \ref{tab:timesB}) of
  the counted epochs.
\item The timing is consistent if $|\Delta t| \leq 3\sqrt{\sigma_{\rm pred}^2 + \sigma_{\rm meas}^2}$. For A this
  tolerance is at least 12.3\,min.
\item The shift of each counted transit alone is reported. It does not enter the outcome.
\end{itemize}

\subsection{Processing variants}

Up to six fixed processing variants are run after the primary analysis and reported next to it: (a) and (b)
detrending windows of 0.5 and 1.0\,d; (c) cadences near SPOC desaturation flags removed within $\pm1$\,h (SPOC primary only); (d) the QLP
light curve (not for a QLP primary; possibly run later, on the release public on 2027 February 28); (e) the uncorrected
simple-aperture flux; (f) the counting rule of an earlier draft, for comparison only. They never
change the primary outcome. They only decide whether an O1 or an O3 is called ``robust'', a word with a narrow meaning:
among the variants (a) to (e) that were run and reach $\Sens \geq 0.95$, none returned O3 against a primary O1, and
none returned O1 or O2 against a primary O3. If none of them reaches 0.95, the qualifier is ``robustness not
testable'', and while variant (d) is pending, ``robust'' and ``robustness not testable'' remain provisional;
a qualifying contradiction already makes the result ``not robust''. The qualifier ``robust'' does not mean that all variants agree
(Appendix~\ref{app:variants}).

\section{Pixel localisation of signal A} \label{sec:localisation}

A dip in a TESS light curve can come from a neighbouring star whose light falls into the same pixels, which are
$21''$ wide. The localisation of a signal is run only if its photometric outcome is O1 or O2. Its position on the sky
is then measured from the Sector 110 target pixel file: each pixel is decorrelated against the spacecraft pointing, a stacked difference image (out of
transit minus in transit) is formed from the counted transits, and the TESS pixel response function (PRF) is fitted to
it \citep{Bell2022}. The offset is measured from the Gaia DR3 position of TIC 4206066 moved to the Sector 110 epoch.

The uncertainty comes from fake transits injected into the real pixels: 300 at the position of the target and 300 at
$5''$ east of it (close to the archival offset of Section~\ref{sec:signals}; this second point was chosen after looking at
the archival data and does not enter the label), and a separate set of 3000 on-target injections for calibration. The measured offset is converted
into a distance $\delta$ in units of the covariance of the on-target cloud, and into an empirical $p$-value,
$p = (1 + m)/(N_{\rm cal} + 1)$, where $N_{\rm cal}$ is the number of successful calibration trials
(2700 to 3000; fewer gives ``not possible'') and $m$ counts those at least as far from the target
\citep{DavisonHinkley1997,North2002,Bates2023}. This $p$ does not assume a Gaussian cloud; it assumes that a real
on-target signal behaves like the injected ones. The label of signal A from Sector 110 alone is
\textbf{L1 compatible} if $p \geq 0.05$, \textbf{L2 tension} if $0.0027 \leq p < 0.05$, and \textbf{L3 off-target} if
$p < 0.0027$. Signal B gets the same measurement without a label, because it has at most three events.

When the localisation is skipped because the photometric outcome is O3 or O4, and in the other cases listed in
Appendices~\ref{app:data}, \ref{app:localisation} and \ref{app:public} (for example a fallback product, failed checks
of the pixel data or of the injections, or a missing or unreadable pixel file), it is reported with the fixed label
``not possible''. A combined position
from Sectors 32, 98 and 110 is also reported; it is descriptive and never overrides the Sector 110 label. With TESS
pixels the position cannot be fixed better than a few arcseconds, so even L1 does not exclude a bound companion or a background star within a few
arcseconds of TIC 4206066 (Section~\ref{sec:limits}).

\section{Limits and what the test cannot show} \label{sec:limits}

\begin{itemize}
\item Whatever the outcome, Sector 110 cannot show whether A or B is a planet on TIC 4206066 rather than an eclipsing
  companion star bound to it and too close to resolve, or a planet transiting such a companion; that needs
  high-resolution imaging and radial velocities. It cannot give
  the mass of either object, which needs radial velocities, or a position better than a few arcseconds.
\item The archival sectors were inspected, the signals were selected after inspection (Section~\ref{sec:signals}), and
  one choice in the ephemeris of A was made after looking at the data (Section~\ref{sec:ephem}).
\item Among the rules set after archival measurements (Appendix~\ref{app:history}), three matter most. The edge zones were introduced after the archival dry runs had
  shown artefacts at data gaps; the two widths tried gave the same result, and the adopted width follows from the
  filter. The 99.9th and 0.1st percentiles were chosen to bring the measured false-alarm level down. This makes detection
  stricter and widens the consistency band slightly. The empirical localisation $p$ replaced a Gaussian $p$ that
  called too many on-target signals ``off target''. The empirical $p$ is more lenient for moderate offsets and gives
  false ``off target'' labels at about the nominal rate. None of these was tuned
  further (Appendix~\ref{app:limits}).
\item The false-detection levels come from trials that share data, so their bounds are approximate and are not a
  guarantee for Sector 110; they are frequencies of the noise model, not probabilities that a detected signal is a
  planet.
\item Signal B will most likely give O4 (Section~\ref{sec:outcomes}); its expected depth rests on one archival transit,
  and its Sector 98 depth is about 1.7 times larger.
\item Edge zones cost data: in a Sector 110 light curve with two or three segments, a transit has roughly a 6 to 10\,\%
  chance to fall in an edge zone and not be tested.
\item The weak 1.76\,d feature seen in Sector 98 (about 130\,ppm) is not part of this test: a period measured in one
  sector cannot predict Sector 110 epochs. Any look at it in Sector 110 is exploratory and will be labelled so, and so
  is any search away from the predicted times.
\end{itemize}
The full list of known limits, stated before the data, is in Appendix~\ref{app:limits}.

\section{The full protocol, which text governs, and reporting} \label{sec:protocol}

\runin{Where the rules are.} Sections~\ref{sec:question} to \ref{sec:protocol} are the main part. Appendices~\ref{app:data}
to \ref{app:codemap} state every rule in full and are part of the pre-registration. The complete pre-registration, with
the appendices, is \code{PREREG\_S110\_FINAL.pdf}; a separate reader-oriented overview is issued as
\code{PREREG\_S110\_SHORT.pdf}. It introduces no additional rules. Appendices A--I give data, processing, thresholds, variants, localisation, procedure, public checks,
fixed wording and commands; J--N give dry runs, limits, history, files and the code's section map.

\runin{Who computes the outcome, and which text governs.} Frozen scripts (deposited with checksums and not changed
afterwards) compute every outcome from valid inputs: \code{s110\_final.py} runs the photometric test,
\code{s110\_localise.py} the pixel localisation, and \code{s110\_summary.py} writes the only sentences allowed to
describe the outcomes. Three rules of precedence apply.
\begin{enumerate}
\item If this text and the frozen code disagree on how an outcome is computed from valid inputs, the code defines the
  outcome. This text defines which runs and inputs are valid (Sections~\ref{sec:ephem} and \ref{sec:times} for the
  ephemeris and predicted times, Section~\ref{sec:stat} for the statistic, and Appendices~\ref{app:data}, \ref{app:processing}, \ref{app:variants},
  \ref{app:localisation}, \ref{app:procedure}, \ref{app:public} and \ref{app:commands}). The code checks only part of
  these rules; the rest are checked through the public record (Appendix~\ref{app:public}), and a few outcomes and
  statements are written by hand from posted files (Appendices~\ref{app:p1}, \ref{app:p5} and \ref{app:p7}). A result obtained in violation of any pre-registered rule
  is invalid and cannot be used to accept or reject either signal.
\item Where Sections~\ref{sec:question} to \ref{sec:protocol} and an appendix differ in any way, or a summary leaves
  something out, the formal rules of the appendix govern.
\item This version 1.1 is the authoritative text of the pre-registration. It replaces the earlier text of 2026 October
  5 (below).
\end{enumerate}

\runin{Order, failures and reporting.} No Sector 110 product is opened before the primary run. The first valid
primary run that writes an output defines the result; it and its checksum are posted before variants or localisation.
Applicable variants run in one continuous session in order (a), (b), (c), (e), (f), then (d) if public; otherwise (d) waits for the
QLP release public on 2027 February 28. Each attempt has an unedited record with command, input and script hashes,
UTC times, exit status and full output. A successful output cannot be replaced by a failure (Appendices F and G).

Before each run, the time-system headers are recorded and checked: TDB, SOLARSYSTEM, reference 2457000, fractional
reference 0 or absent, unit d, and cadence within 1\,\% of 2\,min (SPOC 2-min light curve or pixel file) or 200\,s (QLP/TESS-SPOC). A mismatch gives O4
for the primary, an unavailable variant, or localisation not possible, without conversion or product switching.
An environment or filename refusal is corrected and does not count as an attempt that failed while computing.
A computing failure or interruption before output is repeated once unchanged; a second failure gives O4 (primary), unavailable
(variant), or localisation not possible. A reading adaptation may change only how the file is read; its diff and
exact archival reproduction check must be posted before use. Helpers and individual variants are never adapted.
The binding reading limit is 30 days from release of the affected file, ending at 24:00 UTC; without a qualifying
adaptation it gives O4 (primary) or localisation not possible. A missing pixel file becomes final as not possible on or after 2027 February 28, with a saved search
showing none released by that deadline. The full cases and fixed failure sentences are in Appendix G.

The 30-day primary-run, 7-day posting and 14-day variant/localisation limits are administrative: report delays.
Order, the 2027 February 28 conditions and reading limits are binding. Invalid results are posted as invalid and
cannot accept or reject a signal; if a valid repeat is no longer allowed, the affected outcome is O4 (or localisation
not possible), reason ``pre-registered procedure not followed''. All outcomes, untested transits, variants and
localisations are reported with the fixed wording, including the variant-outcome line wherever ``robust'' is quoted.
The deposit is unchanged after publication; later ideas are separately labelled analyses (Appendices F--H).

\runin{Record status.} Version 1.1 replaces the earlier text deposited on 2026 October 5
(\href{https://doi.org/10.5281/zenodo.23171456}{doi:10.5281/zenodo.23171456}) and deleted the same day; that DOI retains a tombstone.
The public timestamp of this pre-registration is the date of the present deposit.
Unpublished internal drafts v1 and v2 are not part of this deposit and define no outcome.
The scientific choices made after examining archival data are documented in Appendix~\ref{app:history};
editorial changes and the detailed chronology are in \code{CHANGELOG.md}. Frozen code, predictions,
calibrations and archival outputs are unchanged in v1.1.


\begin{acknowledgments}
This work uses data collected by the TESS mission, obtained from the MAST data archive at the Space Telescope Science
Institute (STScI); funding for the TESS mission is provided by NASA's Science Mission Directorate. It uses the TESS-SPOC
(\doi{10.17909/t9-wpz1-8s54}) and QLP (\doi{10.17909/t9-r086-e880}) high-level science products and the TESS PRF
models from MAST. This work has made use of data from the European Space Agency (ESA) mission Gaia
(\url{https://www.cosmos.esa.int/gaia}), processed by the Gaia Data Processing and Analysis Consortium (DPAC,
\url{https://www.cosmos.esa.int/web/gaia/dpac/consortium}). Funding for the DPAC has been provided by national
institutions, in particular the institutions participating in the Gaia Multilateral Agreement.

AI assistance: Claude (Anthropic) assisted code and text drafting; Codex (OpenAI) assisted editing and
verification. I designed the test and am responsible for the protocol. Internal package checks were separate
AI review sessions, not human peer review. Frozen analysis code is unchanged in v1.1.
\end{acknowledgments}

\software{astropy \citepalias{Astropy2013,Astropy2018,Astropy2022}, lightkurve \citepalias{Lightkurve2018}, wotan
\citep{Hippke2019}, TESS\_PRF \citep{Bell2022}}

\appendix

\noindent\textit{Terms used in the appendices.} The ``final re-derivation'' is my complete re-analysis of all archival
data (2026 October 1); its ``items'' 03, 10 and 11 are its transit-time, pixel-localisation and injection analyses, and
the inputs taken from it are in the folder \code{calib/}. ``Audits 1 to 6'' and the ``final independent check'' name the separate AI review sessions recorded in \code{CHANGELOG.md}. ``v1'' and ``v2'' are
the unpublished internal drafts (Section~\ref{sec:protocol}). A ``dry run'' is a run of the frozen code on archival data,
used only to check the code. These names are kept because the files of the deposit use them.

\section{Data products, fallbacks and dates} \label{app:data}

\runin{Primary product.} The test uses the first version of the SPOC 2-min light curve of TIC 4206066 for Sector
110 that is released at MAST, column \code{PDCSAP\_FLUX}, with quality mask 17087 (the default bitmask
\code{TessQualityFlags.DEFAULT\_BITMASK} of lightkurve). Later SPOC reprocessings are reported as additional analyses
and never replace it.

\runin{Fallback products.} A fallback is used only if no SPOC 2-min light curve is public on 2027 February 28. In
that case, on that date and with the version then public:
\begin{enumerate}
\item fallback 1 is the QLP 200-s light curve (column \code{SYS\_RM\_FLUX}, quality mask 1610619071). A corrected QLP
  release counts only if it is public on 2027 February 28;
\item if there is no QLP light curve either, fallback 2 is the TESS-SPOC 200-s full-frame-image light curve
  (\code{PDCSAP\_FLUX}, mask 17087).
\end{enumerate}
A SPOC 2-min product released after a fallback run is reported as an additional analysis, not as the primary.

A fallback run counts only together with the saved result of a MAST search, made on or after 2027 February 28, that
finds no SPOC 2-min light curve of TIC 4206066 for Sector 110. That file and its sha256 checksum are posted with the
fallback result. \code{s110\_summary.py} cannot check this, because the command of fallback 1 is the same as the command
of variant (d) (Appendix~\ref{app:variants}); readers check it. Appendix~\ref{app:p2} fixes the content of this search
record and the consequence of a fallback chosen without it.

\runin{No counted transit.} If the primary product exists but the analysis counts no transit (for example because
\code{PDCSAP\_FLUX} is all NaN), the outcome is O4 for that signal. There is no switch to another product.

\runin{Pixel data.} The localisation uses the first released SPOC 2-min target pixel file of Sector 110. If it does
not exist by 2027 February 28, localisation is reported as not possible.

\runin{Records.} The scripts record the file name and the sha256 checksum of every input.

\section{Light-curve processing} \label{app:processing}

\subsection{Cleaning and detrending}

Cleaning and detrending are unchanged from draft v2, except for the ephemeris used for masking. The light curve is
cleaned with the quality mask of its product. It is then detrended by \code{lcio.detrend} from the frozen helper module
\code{deps/lcio.py}, which uses the biweight filter of wotan \citep{Hippke2019} with a window $W = 0.75$\,d. When the
trend is fitted, both signals are masked over $\pm 1.5\,\Tfourteen$ around their predicted mid-times. Upward outliers
above 4 robust standard deviations ($4 \times 1.4826 \times$ the median absolute deviation) are removed.

\subsection{Segments and edge zones} \label{app:edges}

\code{lcio.detrend} gives wotan the cleaned light curve without the masked windows of both signals ($\pm
1.5\,\Tfourteen$), with a break tolerance of 0.3\,d. Wotan therefore fits a separate segment wherever this series has a
gap longer than 0.3\,d. The edge zones are taken from exactly these segments, for the real detrending and separately
for every fake detrending of the noise model (Appendix~\ref{app:noise}). The first and last $W/2$ of every segment are
edge zones: 0.375\,d (9\,h) for the primary run with $W = 0.75$\,d, and 0.25\,d and 0.5\,d for variants (a) and (b),
whose windows are 0.5\,d and 1.0\,d.

The reason is that the filter is centred and $W$ wide. Within $W/2$ of a segment end it has less than half a window of
data on one side, so its trend there is less reliable. The wotan documentation limits its own edge cut-off to $W/2$
``as this fills the window completely''. A transit whose mask joins a data gap into a break lies between two segments
and is not tested.

\subsection{Counting rule} \label{app:counting}

A predicted transit is counted only if all four conditions hold:
\begin{enumerate}[label=(\roman*)]
\item at least 80\,\% of its $\Tfourteen$ window has cadences after cleaning;
\item it has at least 2 cadences on each side between 0.75 and 3\,$\Tfourteen$ from the predicted mid-time;
\item its whole $\Tfourteen$ window lies inside one segment and outside the edge zones of that segment;
\item in Sector 110, it is one of the frozen epochs of Tables~\ref{tab:timesA} and \ref{tab:timesB}.
\end{enumerate}
A transit that fails any of these conditions is not tested. It is listed with its status and depth, never enters the
statistic, and is never counted as a missing transit.


\subsection{Notes on the expected depths} \label{app:depthnotes}

These notes explain the values of Table~\ref{tab:depths}.
\begin{enumerate}
\item In Sector 32 the two B events fall about 2.0 and 2.3\,h after transits of A. The masks of the two signals then
  join into a hole longer than the 0.3\,d break tolerance, and the detrender fits the two sides separately. By the
  edge-zone rule (Appendix~\ref{app:edges}) these transits are not tested. In Sector 110 the closest pairs (A116 and
  B33, A109 and B31) are 7.4 and 7.7\,h apart, and their masks are separated by 2.1 and 2.4\,h of data. The same
  happens there only if most of that data is lost (about 1.6\,h next to the A mask, about 2.2\,h next to the B mask),
  and then both events of the pair are not tested.
\item The Sector 6 value is measured at 30-minute cadence, where the box statistic is lower than at 2-minute cadence. A
  transit-model check made for draft v2 (kept by me and not part of this deposit, so not verifiable from it; given for information only) gives a ratio of box depth to
  maximum depth of about 0.80 at 30\,min against 0.85 to 0.88 at 2\,min. One of the four counted Sector 6 transits of A
  (epoch $-793$) passes the 80\,\% coverage rule with 0.802, and in Sector 98 epoch 10 also passes with 0.802. The lower box
  depth at 30-minute cadence and these two marginal coverages are accepted as they are.
\end{enumerate}

\section{Noise model, thresholds and measured false-detection levels} \label{app:decision}

\subsection{Noise model} \label{app:noise}

The noise model measures how deep a dip the light curve produces by chance at times where no transit is predicted.
From the cleaned light curve the predicted windows of both signals ($\pm 1.5\,\Tfourteen$) are cut. Fake phase offsets
of the tested signal are then placed between 0.15\,$P$ and 0.85\,$P$, in steps of $\Tfourteen/2$. For each fake phase
the light curve is detrended exactly as in Appendix~\ref{app:processing}, but with the fake ephemeris masked. The
windows at the fake times that pass the counting rule, with the segments and edge zones of that fake detrending, form
the noise pool. Every noise window therefore goes through the same detrending, masking, counting and measurement as a
real window. If the pool has fewer than 100 windows, the outcome is O4.

As a supplementary check, the whole set of counted windows at each fake phase gives a full-set statistic $d_\phi$, and
the number of fake phases with $d_\phi \geq d$ is reported. This number is not part of the outcome: it is in-sample, its
sets have their own number of windows, and, because the phases overlap, the count can only distinguish levels of about
5\,\%.

\subsection{Thresholds} \label{app:thresholds}

For each signal, with $n$ counted transits and statistic $d$:
\begin{itemize}
\item the noise pool is bootstrapped 20\,000 times to give the distribution of the mean of $n$ windows;
\item $\ddet$ is the 99.9th percentile of this distribution;
\item $\dlo$ is the 0.1st percentile of (noise $+$ a depth drawn from $N(D, D_{\rm err})$);
\item $\dhi$ is the 99.9th percentile of (noise $+$ a depth drawn from $N(D_{98}, D_{98,\rm err})$);
\item the sensitivity $\Sens = P(\mathrm{detect} \mid \mathrm{real})$ is the fraction of (noise $+$ $N(D, D_{\rm err})$)
  above $\ddet$.
\end{itemize}
All four quantities are computed for 20 bootstrap seeds, and their medians are used. The 99.9th percentile of 20\,000
draws rests on about 20 draws, so $\ddet$ moves between seeds (for B by up to about 10\,\%). The seed spread is
reported and plays no role in the outcome. A value of $d$ inside the spread is not described as ``marginal''.

\subsection{Measured false-detection levels} \label{app:fa}

The bootstrap of a finite pool of overlapping windows underestimates the tails of the noise, so the 99.9th percentile
is not the false-alarm level of the test. The levels that actually apply were therefore measured with a separate set of fake
ephemerides on the archival sectors: the true transits were removed, the fake periods were arbitrary, each fake phase
was one trial, and the thresholds were computed from the pool without the overlapping phases. The levels were measured
per product and per number of counted transits $n$ (Table~\ref{tab:fa}). The level that applies to Sector 110 is the
one of its product (SPOC 2-min for the primary run) and of the $n$ actually counted.

\begingroup\footnotesize
\begin{longtable}{>{\raggedright\arraybackslash}p{0.05\textwidth}>{\raggedright\arraybackslash}p{0.162\textwidth}>{\raggedright\arraybackslash}p{0.162\textwidth}>{\raggedright\arraybackslash}p{0.162\textwidth}>{\raggedright\arraybackslash}p{0.162\textwidth}>{\raggedright\arraybackslash}p{0.162\textwidth}}
\caption{Measured false-detection levels on archival data, from \code{calib/false\_alarm\_levels.json}: exceedances
$k$ out of $N$ trials. The first row pools all trials of a product and signal, with the 95\,\% interval from a bootstrap
over the trial ephemerides; the other rows give exactly $n$ counted transits. ``Bound'': one-sided 95\,\% Clopper-Pearson
upper bound for independent trials (see text). Empty: no trial with this $n$. The five columns contain all
17\,835 trials. The TESS-SPOC 30-min trials (Sector 6) are not quoted by \code{s110\_summary.py}: the Sector 110
TESS-SPOC product has 200-s cadence, and for it the SPOC 2-min levels are quoted.\label{tab:fa}}\\
\toprule
$n$ & SPOC 2-min (primary), A & SPOC 2-min (primary), B & QLP 200-s (fallback 1), A & QLP 200-s (fallback 1), B & TESS-SPOC 30-min (Sector 6), A \\
\midrule
\endfirsthead
\toprule
$n$ & SPOC 2-min (primary), A & SPOC 2-min (primary), B & QLP 200-s (fallback 1), A & QLP 200-s (fallback 1), B & TESS-SPOC 30-min (Sector 6), A \\
\midrule
\endhead
all & 0.47\,\% (26/5514; interval 0.27 to 0.70; bound 0.65\,\%) & 0.49\,\% (51/10312; interval 0.35 to 0.64; bound 0.62\,\%) & 0/614 (bound 0.49\,\%) & 0.64\,\% (5/777; interval 0.29 to 0.94; bound 1.35\,\%) & 0/618 (bound 0.48\,\%) \\
1 & 0/23 (bound 12.21\,\%) & 20/3054 (0.65\,\%, bound 0.95\,\%) & 0/1 (bound 95.00\,\%) & 2/188 (1.06\,\%, bound 3.31\,\%) & 0/11 (bound 23.84\,\%) \\
2 & 1/68 (1.47\,\%, bound 6.79\,\%) & 31/6040 (0.51\,\%, bound 0.69\,\%) & 0/6 (bound 39.30\,\%) & 2/325 (0.62\,\%, bound 1.92\,\%) & 0/51 (bound 5.70\,\%) \\
3 & 7/642 (1.09\,\%, bound 2.04\,\%) & 0/1077 (bound 0.28\,\%) & 0/33 (bound 8.68\,\%) & 1/202 (0.50\,\%, bound 2.33\,\%) & 0/140 (bound 2.12\,\%) \\
4 & 6/1538 (0.39\,\%, bound 0.77\,\%) & 0/141 (bound 2.10\,\%) & 0/136 (bound 2.18\,\%) & 0/62 (bound 4.72\,\%) & 0/175 (bound 1.70\,\%) \\
5 & 1/1203 (0.08\,\%, bound 0.39\,\%) &  & 0/197 (bound 1.51\,\%) &  & 0/118 (bound 2.51\,\%) \\
6 & 5/963 (0.52\,\%, bound 1.09\,\%) &  & 0/102 (bound 2.89\,\%) &  & 0/71 (bound 4.13\,\%) \\
7 & 6/907 (0.66\,\%, bound 1.30\,\%) &  & 0/66 (bound 4.44\,\%) &  & 0/51 (bound 5.70\,\%) \\
8 & 0/146 (bound 2.03\,\%) &  & 0/35 (bound 8.20\,\%) &  & 0/1 (bound 95.00\,\%) \\
9 & 0/24 (bound 11.73\,\%) &  & 0/32 (bound 8.94\,\%) &  &  \\
10 &  &  & 0/4 (bound 52.71\,\%) &  &  \\
11 &  &  & 0/2 (bound 77.64\,\%) &  &  \\
\bottomrule
\end{longtable}
\endgroup

Every trial set that was run is included: the three-product sets, two later SPOC 2-min sets, and the SPOC 2-min set run
by audit 3 (\code{build/false\_alarm\_check.py}, modes A, B, A\_S32, B\_S32, A\_S32\_a3 and B\_S32\_a3). Three of audit
3's ephemerides were recomputed by the trial script of the deposit and agree exactly. Since audit 6 every trial uses
exactly the thresholds of the Sector 110 run (\code{s110\_final.thresholds}, the median of 20 seeds). Before that, the
trial script used one seed, which changed 14 of the 17\,835 trial classifications (4 more and 10 fewer exceedances; all
trial depths identical).

``Bound'' is the one-sided 95\,\% Clopper-Pearson upper bound for independent trials \citep{ClopperPearson1934}. The
trials are not independent: neighbouring phases share data, and all trials of one product share one light curve. The
bound is therefore an approximation whose coverage is not established. It is not a guaranteed maximum, and neither the
bound nor the interval is a guarantee for Sector 110. These levels are frequencies of the noise model per trial, not
probabilities that a detected signal is a planet. At a given $n$, the level of B is not always higher than that of A. For an $n$ with few
trials the point value can be 0 while its bound is several times the overall level (A, $n = 8$: 0 of 146 trials, bound
2.03\,\%). The sentence of Appendix~\ref{app:wording} prints both, and the bound is the number to read. A grouping of $n$
that was tried after seeing the per-$n$ counts is not used. \code{s110\_summary.py} quotes, for the product of the run,
the overall level and the level for exactly the $n$ counted, each with its binomial bound and the remark on dependent
trials.

A check of the lower tail gave the matching numbers for real signals: a signal at the expected depth fell below $\dlo$
in 0.10\,\% (A) and 0.17\,\% (B) of the A\_S32 and B\_S32 trials. For fallback 2 (TESS-SPOC 200-s) no trials exist,
and the SPOC 2-min levels are quoted with that remark.

The trial files are \code{calib/false\_alarm\_*.json}, made by \code{build/false\_alarm\_check.py}; the levels are in
\code{calib/false\_alarm\_levels.json}, built by \code{build/false\_alarm\_levels.py}. The intervals come from a
bootstrap over the trial ephemerides. Because all trials of one product share one light curve, these intervals are
optimistic. The percentiles were set once (99.9 and 0.1) and are not tuned further. The levels are reported as they are,
and \code{s110\_summary.py} quotes the level of the product and of the $n$ of the actual run.

The transit times are fixed in advance and nothing is searched, so these are fixed-epoch levels, not search-wide
false-alarm rates. There is no minimum number of transits: the thresholds and the sensitivity are computed for the $n$
actually counted. An exclusion (O3) needs $\Sens \geq 0.95$ for that $n$. Whether this is reached depends on the
expected depth, the noise and $n$, not on a fixed number of transits; in the archival trials some sets reach it with
$n = 2$. A detection is possible for any $n$, and its false-detection level is the one of Table~\ref{tab:fa} for that
$n$.

\section{Processing variants and the robustness qualifier} \label{app:variants}

\subsection{Variants}

After the primary run, fixed alternative processings (variants) are run on the same product and reported next to the
primary result. They use the expected depths of the primary run; they are not recalibrated per variant. The variants
are:
\begin{enumerate}[label=(\alph*)]
\item detrending window 0.5\,d;
\item detrending window 1.0\,d;
\item cadences within $\pm 1$\,h of SPOC desaturation flags removed (only when the primary product is a SPOC product);
\item the QLP light curve with mask 1610619071: the first QLP release that is public when the other variants are run,
  or, if none is public then, the release that is public on 2027 February 28. It is run once and not re-run on later
  releases. It is not run when the primary product is itself QLP;
\item the \code{SAP\_FLUX} of the same product. This flux is not corrected for light from other stars, so variant (e)
  is diluted and leans toward ``shallower'';
\item for comparison only, outside the robustness rule: the counting rule of draft v2 (80\,\% coverage only, no edge
  zones, so every transit with 80\,\% coverage is counted).
\end{enumerate}

Variants are never edited. A variant that crashes or cannot be run is ``not available'' and is ignored. ``Not
available'' means one of two things: the registered command of the variant, run unchanged on the registered file,
stopped with an error while reading or computing (after the one repeat allowed by Appendix~\ref{app:procedure}, point
4); or its input file fails the time-system check of Appendix~\ref{app:p3}. Its attempt record
(Appendix~\ref{app:p4}), which contains the unedited error output or the printed header line, is saved and posted. A
refusal by the environment check or the file-name check is not ``not available'': the environment or the file name is
corrected and the variant is run. For variant (d), ``not yet public'' is handled by the robustness rule below.
\code{s110\_summary.py} copies the error output but cannot judge it. Because data and code are public, anyone can rerun
a variant declared not available. Appendix~\ref{app:p4} fixes what the saved record of an attempt must contain and when
a ``not available'' is invalid.

\subsection{Robustness qualifier} \label{app:robust}

Variants never change any primary outcome. They only decide whether an O1 or an O3 carries the qualifier ``robust''.
Only variants (a) to (e) with $\Sens \geq 0.95$ (qualifying variants) are considered.
\begin{itemize}
\item If no variant among (a) to (e) reaches $\Sens \geq 0.95$, the qualifier is ``robustness not testable''.
\item If at least one variant reaches $\Sens \geq 0.95$: a primary O1 is called robust only if no qualifying variant returns
  O3, and a primary O3 is called robust only if no qualifying variant returns O1 or O2.
\item O2 and O4 carry no robustness qualifier.
\end{itemize}

If variant (d) is not yet public when the other variants are run, ``robust'' and ``robustness not testable'' are
reported as provisional. A qualifying contradiction already gives ``not robust'', without a provisional qualifier:
a later variant cannot remove that contradiction. A provisional qualifier becomes final on 2027 February 28, with
(d) if (d) is public by then. ``Not yet public'' keeps the qualifier provisional; it
never makes it final. A provisional qualifier becomes final only in one of three cases: (d) has been run; a public QLP release
was attempted and the run failed; or a MAST search made on or after 2027 February 28 finds no QLP Sector 110 release.
\code{s110\_summary.py} accepts the last case only on or after that date (audit 6).

The qualifier is computed by \code{s110\_summary.py}, not by hand. The only exception is the cases of
Appendix~\ref{app:p5}, where the summary script cannot run; there the qualifier is written from the posted variant
results by exactly this rule. ``Robust'' has exactly the meaning given here, which is narrower than ``all variants
agree''. Appendix~\ref{app:p7} fixes this meaning and the line of variant outcomes that is always posted with it.

\subsection{Exploratory analyses}

Aperture photometry of my own and any search away from the predicted times (shifted transits, other periods) are
exploratory and cannot be cited for or against any outcome.

\section{Pixel localisation: full rules} \label{app:localisation}

\subsection{When it is run, and with which code}

The localisation of a signal is run only if its photometric outcome is O1 or O2. Otherwise it is reported as not
possible.

The code is \code{s110\_localise.py}. It runs the headline recipe of the final re-derivation (items 10 and 11) through
the library \code{deps/pixloc\_lib.py}, which is the re-derivation's code. On 2026 October 5 only the comments of this
library and its default file paths, which \code{s110\_localise.py} replaces at run time, were made neutral. On the
Sector 32 pixels with the re-derivation's ephemeris the code reproduces the published Sector 32 result exactly,
including all 600 injection positions (\code{dryrun/loc\_s32\_items10\_fidelity.json}).

\subsection{Recipe}

\begin{itemize}
\item Input: the Sector 110 SPOC 2-min target pixel file. Cadences with quality bits 1610619071 are removed.
\item Each pixel is decorrelated against the spacecraft pointing columns \code{POS\_CORR1} and \code{POS\_CORR2}, with
  a linear fit made away from the localised signal's own predicted transits, as in item 10.
\item A stacked difference image is formed: the mean of the windows from 1.5 to 2.5\,$\Tfourteen$ before and after each
  transit, minus the in-transit mean.
\item A PRF is fitted to the difference image, weighted by pixel errors based on \code{FLUX\_ERR} (TESS\_PRF 0.1.3
  \citep{Bell2022}, MAST PRF files \code{start\_s0004}, camera and CCD taken from the file).
\item Offsets are measured relative to the Gaia DR3 position, moved with its proper motion to the Sector 110 epoch.
\item Transits used: the counted transits of the primary photometric run. The pixel code also requires both
  out-of-transit windows of a transit; any transit it drops is listed.
\end{itemize}

\subsection{Error clouds from injections} \label{app:clouds}

The uncertainty of the position comes from fake transits injected into the real pixels, using the real stellar image
as template:
\begin{itemize}
\item 300 injections at the position of the target (seed 1001) and 300 at $+5''$ east (seed 2002);
\item for calibration, 3000 further on-target injections with the same recipe (seed 4004).
\end{itemize}
Every fake stack has exactly as many transits as the real stack. A candidate fake transit must lie more than
$2.5\,\Tfourteen + \Tfourteen^{\rm other}/2$ from the other signal, pass the pixel usability rule, have its $\Tfourteen$
window inside one segment of the primary photometric run at least $W/2$ from its ends (the photometric edge-zone rule;
the segments are read from the photometry output file), and lie more than 3\,$\Tfourteen$ from every fake already chosen
in that trial. For each trial a phase is drawn and its candidates are taken. While fewer candidates than needed are
collected, further phases are drawn and their candidates added; no phase is rejected, and after at most 100 phases the
trial fails. The needed number is then chosen at random from the candidates. Each cloud records the number of distinct
$\Tfourteen$-wide phase slots of its fakes. Fewer than 15 slots in any cloud make the localisation ``not possible''.
The injected depth is the pixel-aperture box depth of the real transits, as in item 11.

\subsection{What is reported}

\begin{itemize}
\item The offset (East, North) of the dip from the target.
\item Its distance $\delta$ from the target position $(0, 0)$, in units of the covariance of the 300-trial on-target
  cloud (a Mahalanobis distance; called $d$ in the output files and in the earlier drafts).
\item The empirical $p = (1 + m)/(N_{\rm cal} + 1)$, where $N_{\rm cal}$ is the number of successful calibration
  trials (2700 to 3000) and $m$ counts those whose distance $\delta_i$, measured in the same way, is at least
  $\delta$. When all trials succeed, the denominator is $3000+1$. This states the existing frozen-code
  calculation; fewer than 2700 successful trials gives ``not possible''. This is a Monte Carlo $p$-value with a separate
  calibration set \citep{DavisonHinkley1997,North2002,Bates2023}.
\item $L(5''\,{\rm E})/L({\rm target})$, the ratio of the Gaussian densities of the two clouds at the observed position.
  It is a ratio, not a probability, and the $5''$ east hypothesis was defined after looking at archival data.
\end{itemize}
The label uses the empirical $p$. It does not assume that the cloud is Gaussian; it assumes only that a real on-target
signal behaves like the injected ones. The Gaussian tail $\exp(-\delta^2/2)$ and a $p$ measured from the mean of the
injection cloud (``bias-corrected'') are reported for information and never change the label. The coverage of the
empirical $p$ was checked on archival and synthetic pixels (Appendices~\ref{app:validation} and \ref{app:limits}).

\subsection{Labels for signal A}

The label of signal A from Sector 110 alone is:
\begin{itemize}
\item \textbf{L1 compatible} if $p \geq 0.05$;
\item \textbf{L2 tension} if $0.0027 \leq p < 0.05$;
\item \textbf{L3 off-target} if $p < 0.0027$.
\end{itemize}

The localisation is \textbf{not possible} in any of these cases: the pointing columns are missing, all zero or
non-finite on more than 50\,\% of cadences (as coded: \code{POS\_CORR1} or \code{POS\_CORR2} is missing, or fewer than
half of the cadences have both values finite and not both zero); no transit passes the pixel rule; fewer than 90\,\% of the
trials succeed in any cloud (270 of 300, or 2700 of 3000); or $\delta$ is not finite.

\subsection{Checks before a labelled run}

The script refuses a labelled run unless all of the following hold: the installed package versions equal
\code{requirements.txt}; the photometry file comes from the primary settings, and its sha256 equals the publicly posted
primary checksum (Appendix~\ref{app:procedure}); its input is the SPOC 2-min light curve of Sector 110; the pixel file is
the matching SPOC 2-min target pixel file of TIC 4206066; the PRF folder holds exactly the files listed in
\code{PRF\_SHA256SUMS}, with those checksums; and the registered switches are used. The environment is checked before
any file is read, and every other check comes before any rule that writes a result.

Only after these checks: if the pixel file is not from camera 1, CCD 2 (the registered PRF files), localisation is not
possible. If the primary product is a fallback, localisation is not possible. The switch \code{--dry-run} is refused on
Sector 110 files and never produces a label. The first labelled run of each signal that writes its output file defines
the result; that file and its sha256 are posted (Appendix~\ref{app:procedure}). Any other localisation of Sector 110 is
exploratory and cannot be cited for or against the label.

\subsection{Combined position and signal B}

The positions of Sectors 32, 98 and 110 are also combined: an inverse-covariance mean of the three positions, using the
archival positions and covariances of the re-derivation headline (\code{calib/localisation\_archival.json}; Sector 32
decorrelated and Sector 98 ``hard'' mask, main recipe), and the product of the three likelihood ratios. The combined
value is descriptive and never overrides the Sector 110 label. The archival values were measured with the
re-derivation's older ephemeris (about 3\,min different), which moves no window by more than 3\,\%.

Signal B gets the same measurement without a label, because it has at most three events, with injections kept clear
of signal A.

\section{Procedure and commitments} \label{app:procedure}

\begin{enumerate}
\item \textbf{Order.} No Sector 110 product of any kind (light curves, pixels, full-frame images, SPOC data-validation
  reports, threshold-crossing-event or TESS Object of Interest lists for this star) is opened before the primary run.
  The primary run is made within 30 days of the release of the primary product (Appendix~\ref{app:data}). A later run is
  still the primary run, and its delay is reported.
\item \textbf{Refusals.} Every registered run (photometry, variants, localisation) refuses to start, before reading any
  data, in three cases: the installed package versions, the machine type or the operating system differ from
  \code{requirements.txt}; the input is not Sector 110 of TIC 4206066; or the file name is not the registered name of
  the product. If MAST distributes a product under another name, the file is renamed to the registered pattern of
  \code{s110\_final.py}, and the original name and its sha256 are recorded. A refused run writes nothing and is not a
  run. The exact commands, run from the root of the deposit, are listed in
  Appendix~\ref{app:commands}.
\item \textbf{First run and posting.} The first primary run that writes its output file is the primary run. That file
  (which contains the input, code and dependency checksums and the package versions) and its sha256 are posted as a new
  version of this Zenodo record within 7 days of the run. Only then are the variants and the localisation run, within
  14 days of posting the primary result. Their output files and the summary file are posted in the same way within 7
  days of being made.
\item \textbf{Failures and repeats.} If a script writes its output file, that file defines its outcomes, even if a later
  step fails. A run that ends before writing its output for any reason other than a refusal or a failure to read the
  file (for example an interruption) is repeated once, unchanged. If it fails again, the outcome is O4 for both signals
  (photometry) or ``localisation not possible''. For a reading failure the unedited error message of the script is saved
  and its checksum recorded. A code edit is allowed only to read the file (column or extension names, data type), and
  the edited script keeps its registered file name (Appendix~\ref{app:p6}). It may not change any constant, window,
  mask, seed, threshold or rule. Before use, the edited script must reproduce exactly the \code{results} block of every
  photometric \code{dryrun/*.json}, and the \code{stats} block and every \code{injections.sets.*.trials} list of every
  localisation \code{dryrun/loc\_*.json} (run times differ and are not compared). The diff and this check are posted, as
  in point 3, before the Sector 110 run. Variants are never edited. Any other failure of the localisation gives
  ``localisation not possible''.
\item \textbf{Full reporting.} All outcomes are reported, whatever they are, for A and B, with all variants, the
  transits that were not tested and the reasons, and the localisation.
\item \textbf{No changes.} Nothing in this deposit is changed after publication. Any later idea is a new, separately
  labelled analysis.
\item \textbf{Public checks.} Readers can check every posted output file: its \code{script\_sha256},
  \code{deps\_sha256}, \code{ephemeris\_sha256} and \code{expected\_depths\_sha256} must equal the checksums in
  \code{SHA256SUMS} of this deposit (or, for a script adapted under point 4 to read the file, the checksum of the
  adapted script whose diff and check were posted; Appendix~\ref{app:p6}). The scripts compare these checksums only
  with the files on the computer that runs them, and \code{s110\_summary.py} takes the date of its 2027 February 28
  rule from that computer's clock. The public comparison is therefore the check that counts. A fallback primary needs
  the MAST search of Appendix~\ref{app:data}, and a variant declared not available needs its error output
  (Appendix~\ref{app:variants}). Appendix~\ref{app:public} lists every rule that only this public comparison checks.
\end{enumerate}

\section{Rules checked through the public record} \label{app:public}

The scripts check the software environment, the sector, the star, the file name, the registered arguments and the
checksums that every result carries (Appendices~\ref{app:localisation} and \ref{app:procedure}). They do \emph{not}
check the rules of this section. These rules are part of the pre-registration exactly like the coded ones. They are
checked through the public record: every file named below is posted with its sha256 as in
Appendix~\ref{app:procedure}, point 3, and because the data and the code are public, anyone can repeat each step.
Nothing in this section claims that the code prevents a violation.

\subsection{Invalid results} \label{app:p1}

\begin{itemize}
\item A result (a photometric, variant or localisation output file, a summary file, or a sentence) obtained in
  violation of a pre-registered rule on the ephemeris and the predicted times (Sections~\ref{sec:ephem} and
  \ref{sec:times}), the statistic (Section~\ref{sec:stat}), or of a rule of Appendices~\ref{app:data}, \ref{app:processing}, \ref{app:variants},
  \ref{app:localisation}, \ref{app:procedure}, \ref{app:public} or \ref{app:commands}, is \emph{invalid}. An invalid result is not an
  outcome of this test. It cannot be used to accept or to reject signal A or signal B, and it is not described with the
  sentences of Appendix~\ref{app:wording}. It is still posted, labelled ``invalid'' with the rule it breaks
  (Appendix~\ref{app:procedure}, point 5).
\item That a script accepted a run, or that \code{s110\_summary.py} wrote a sentence, does not by itself make the result
  valid.
\item If a violation is found after posting, by me or by anyone, a new version of this record declares the affected
  results invalid and names the rule. Where the rules still allow it, the step is repeated as registered (for example
  the primary run on the product that Appendix~\ref{app:data} requires). An invalid run never counts as the first run of
  Appendix~\ref{app:procedure}, point 3; the repeated run defines the outcome, and the delay and the earlier invalid run
  are reported. Where the rules no longer allow a repeat, the outcome of each affected signal is O4 with the reason
  ``pre-registered procedure not followed'', or, for the localisation, ``localisation not possible'' with that reason.
  This is the case when the date 2027 February 28 or a reading limit of Appendix~\ref{app:p5} has passed, or when a rule
  of order of Appendix~\ref{app:procedure}, points 1 and 3, was broken, which a repeat cannot repair.
\item The time limits of Appendix~\ref{app:procedure}, points 1 and 3, are administrative: the primary run within 30 days
  of the release of the primary product; results posted within 7 days; variants and localisation run within 14 days of
  posting the primary result. No outcome, threshold, label or sentence depends on them, and none of the scripts reads
  them (the only date any script compares is 2027 February 28, in \code{s110\_summary.py}). A run or a posting after
  such a limit is reported with its delay; it is not invalid for that reason alone. Binding, as before, are every rule
  of order (no Sector 110 product opened before the primary run, Appendix~\ref{app:procedure}, point 1; variants and
  localisation only after the primary result is posted, point 3), the date 2027 February 28
  (Appendices~\ref{app:data}, \ref{app:variants}, \ref{app:p2} and \ref{app:p5}), and the reading limits of
  Appendix~\ref{app:p5}, which decide an outcome.
\item Two different 30-day limits exist, and they are never interchangeable.
  \begin{itemize}
  \item The \emph{run limit} (Appendix~\ref{app:procedure}, point 1): the primary run is made within 30 days of the
    release of the primary product. It is administrative, as stated above: a later primary run is still the primary
    run, with its delay reported.
  \item The \emph{reading limits} (Appendix~\ref{app:p5}): if a file cannot be read by the registered script, an adapted
    reader that passes the check of Appendix~\ref{app:procedure}, point 4, must be posted within 30 days of the release
    of that file at MAST. The file is the light curve for the photometric run and the target pixel file for the
    localisation. The release date is the one listed by the product search record of Appendix~\ref{app:p2}, or, where
    MAST lists none, the date of the first product search that lists the file. The limit ends at 24:00 UTC of the 30th
    day. The reading limits are binding: when one has passed without such a posting, the outcome of the corresponding
    row of Table~\ref{tab:failures} applies.
  \end{itemize}
  Meeting the run limit never meets a reading limit, and the administrative status of the run limit never applies to a
  reading limit. Each of the two reading limits counts only from the release of its own file.
\end{itemize}

\subsection{Choice of the product (fallbacks)} \label{app:p2}

\begin{itemize}
\item Before the primary run a product search record is saved. It contains the exact MAST query (the astroquery call
  with all its arguments, or the MAST Portal search with its URL), the UTC date and time of the search, and the
  complete, unfiltered list of Sector 110 products of TIC 4206066 that the search returns (file names, product types,
  and the release date or data release number where MAST gives one). It names the light curve and the target pixel
  file selected, with their sha256 checksums, and, for a renamed file (Appendix~\ref{app:procedure}, point 2), the
  original name. Listing, downloading and computing checksums of these files is not an opening of the product in the
  sense of Appendix~\ref{app:procedure}, point 1.
\item The primary run is valid only on the first released SPOC 2-min light curve listed by this record. A fallback run
  is valid only if the record was made on or after 2027 February 28 and lists no SPOC 2-min light curve (fallback 1,
  QLP), or lists neither a SPOC 2-min nor a QLP light curve (fallback 2, TESS-SPOC). A fallback run made before
  2027 February 28, without this record, or while the record lists a product of higher priority is invalid, whatever
  its outcome.
\item The code does not check any of this: \code{s110\_final.py} accepts a QLP or TESS-SPOC file with a registered file
  name as a registered run on any date, and \code{s110\_summary.py} writes sentences for it.
\item Variant (d) uses a record of the same kind for its QLP file (Appendix~\ref{app:variants}). The record for the
  primary product is posted with the primary result. The record for (d) is posted with the result of (d), or, if (d) is
  not run, with the summary that declares it not available.
\end{itemize}

\subsection{Time system of every input file} \label{app:p3}

\begin{itemize}
\item Every analysed file must declare in its first extension (HDU 1) the time system of all archival inputs of this
  test, which was checked on the archival SPOC 2-min, TESS-SPOC and QLP files of Sectors 6, 32 and 98 on 2026 October 5:
  \code{TIMESYS} $=$ TDB, \code{TIMEREF} $=$ SOLARSYSTEM, \code{BJDREFI} $=$ 2457000, \code{BJDREFF} $=$ 0 or absent (QLP
  files do not carry it), \code{TIMEUNIT} $=$ d, and \code{TIMEDEL} equal, within 1\,\%, to 2\,min for a SPOC 2-min light
  curve or target pixel file and to 200\,s for a QLP or TESS-SPOC 200-s light curve.
\item As the first step of every run, before the registered command, these keys are printed without reading any data
  column, with the command given in Appendix~\ref{app:commands}. The printed line is kept in the attempt record
  (Appendix~\ref{app:p4}). Reading these keys is part of the run, not an opening of the product in the sense of
  Appendix~\ref{app:procedure}, point 1.
\item If any key differs, the file is not analysed, no conversion is made, and there is no switch to another product.
  For the primary file the outcome is O4 for both signals, with the reason ``time system of the product not as
  registered''. For a variant file the variant is not available, with the printed line as its record. For the pixel
  file the localisation is not possible, with the reason ``time system of the pixel file not as registered''.
\item A run on a file whose keys were not printed and kept, or differ, is invalid. \code{s110\_final.py} and
  \code{s110\_localise.py} do not read these keys.
\end{itemize}

\subsection{Attempt records and ``not available''} \label{app:p4}

\begin{itemize}
\item Every attempt of every registered command (primary or fallback, each variant, each localisation, the summary)
  has a plain-text attempt record, written while the run happens and not edited afterwards. It contains the exact
  command line; the name and sha256 of every input file; the sha256 of the script run; for a run on a FITS file, the
  printed line of Appendix~\ref{app:p3}; the UTC start and end time; the exit status; and the complete standard output and
  standard error. The record is written by the macOS \code{script} utility in UTC, for example
  \code{TZ=UTC script -a RECORD <command>}, which saves the command, its start and end time, all its output and its exit
  status. The checksum lines of the inputs and of the script, and the line of Appendix~\ref{app:p3}, are written into
  the same record in the same way just before the registered command. There is one record per attempt, also for
  successful attempts and for refusals. All records are posted with the output files they belong to
  (Appendix~\ref{app:procedure}, point 3), so the number and the order of the attempts are public.
\item The variants are attempted in one continuous session, in the order (a), (b), (c), (e), (f), and then (d) if a QLP
  Sector 110 release is public when the session starts. Each is attempted once and, after a failure, repeated once
  (Appendix~\ref{app:procedure}, point 4). The order is binding. A session that runs past midnight UTC continues in the
  same order; its start and end are in the records, and nothing else follows from the date. If no QLP Sector 110
  release is public when the session starts, (d) is attempted later, in the same way and with its own records, on the
  QLP release that is public on 2027 February 28 (Appendix~\ref{app:variants}), on or after that date. If MAST lists no
  QLP Sector 110 release on 2027 February 28, (d) is final as ``not-public-on-deadline''.
\item A variant is ``not available'' only if its two records show the registered command on the registered input (for
  (a) to (c), (e) and (f), the primary file, with the same sha256), each ending with a non-zero exit status after an
  error while reading or computing. For a reading failure, which Appendix~\ref{app:procedure}, point 4, does not repeat,
  one such record is enough. A variant is also not available if its file fails the check of Appendix~\ref{app:p3}. The
  file given to \code{--not-available} is the record of the last attempt (or the printed line of
  Appendix~\ref{app:p3}).
\item A variant whose registered run wrote its output file can never be declared not available; that file is its
  result. A refusal is not ``not available'' (Appendix~\ref{app:variants}).
\item A summary that declares a variant not available with any other file (free text, a refusal, an edited or
  shortened output, the record of another command) is invalid. \code{s110\_summary.py} accepts any non-empty file there
  and does not check its content.
\item Variant (d) is final as ``crashed'' with the two records of the attempt on a QLP release that was public at the
  time. It is final as ``not-public-on-deadline'' with the product search record of Appendix~\ref{app:p2}, made on or
  after 2027 February 28, that lists no QLP Sector 110 light curve. A file for (d) that fails the check of
  Appendix~\ref{app:p3} is declared final with \code{--d-final-reason crashed} and the printed line of
  Appendix~\ref{app:p3} as the file.
\item When a variant among (a) to (e) is not available, the robustness qualifier is computed without it. The line of
  Appendix~\ref{app:p7} shows it as ``not available'', and the qualifier is never described as covering it.
\end{itemize}

\subsection{Registered failures that leave no output file} \label{app:p5}

Some registered outcomes are reached without an output file, so \code{s110\_summary.py} cannot write their sentences
(it needs the primary result and, for a SPOC primary, the localisation result of each signal with O1 or O2). For these
cases only, the sentence is written from the posted files exactly by the tables of Appendix~\ref{app:wording}; only the
reason comes from Table~\ref{tab:failures}.

\begin{longtable}{>{\raggedright\arraybackslash}p{0.34\textwidth}>{\raggedright\arraybackslash}p{0.16\textwidth}>{\raggedright\arraybackslash}p{0.14\textwidth}>{\raggedright\arraybackslash}p{0.26\textwidth}}
\caption{Registered failures that leave no output file.\label{tab:failures}}\\
\toprule
Case & Posted record & Outcome & Fixed reason R \\
\midrule
\endfirsthead
\toprule
Case & Posted record & Outcome & Fixed reason R \\
\midrule
\endhead
the primary photometric run ends twice without writing its output (not a refusal, not a reading failure;
  Appendix~\ref{app:procedure}, point 4) & the two attempt records & O4 for A and B & ``the registered photometric run
  failed twice without writing a result'' \\
the primary file cannot be read, and no adapted reader (Appendix~\ref{app:p6}) that passes the check of
  Appendix~\ref{app:procedure}, point 4, is posted within 30 days of the release of the primary product & the attempt
  record with the error output & O4 for A and B & ``the registered product could not be read'' \\
the primary file fails the check of Appendix~\ref{app:p3} & the printed line & O4 for A and B & ``time system of the
  product not as registered'' \\
localisation is needed, and a product search of 2027 February 28 or later lists no SPOC 2-min target pixel file of
  Sector 110 & that search record & localisation not possible & ``no SPOC 2-min target pixel file released by
  2027-02-28'' \\
the localisation run ends twice without writing its output & the two attempt records & localisation not possible &
  ``the registered localisation run failed twice without writing a result'' \\
the pixel file cannot be read, and no adapted reader is posted within 30 days of its release & the attempt record
  with the error output & localisation not possible & ``the target pixel file could not be read'' \\
the pixel file fails the check of Appendix~\ref{app:p3} & the printed line & localisation not possible & ``time system
  of the pixel file not as registered'' \\
\bottomrule
\end{longtable}

\begin{itemize}
\item Photometric failure: for each signal X, ``Sector 110 cannot decide on signal X (reason: R).''
\item Localisation failure: the photometric part and the robustness part are written exactly as
  Appendix~\ref{app:wording} defines them, from the posted primary and variant results and
  \code{calib/false\_alarm\_levels.json}. The localisation part is ``; its position could not be measured (R).'' for A
  and ``; B's localisation was not possible (R).'' for B, written like the localisation parts of
  Appendix~\ref{app:wording}, without the word ``reason''. The sentence of the other signal is written in the same way
  from its own files.
\item These sentences are not produced or checked by the code; anyone can rebuild them from the tables and the posted
  files.
\item A pixel file that is missing before 2027 February 28 is not a final failure: the sentences wait until the file is
  released or the date passes. The primary result is posted in any case (Appendix~\ref{app:procedure}, point 3).
\item A failure never replaces a result: if any registered attempt of a step wrote its output file, that file defines
  the result (Appendix~\ref{app:procedure}, point 4), and a failure sentence for that step is invalid.
\end{itemize}

\subsection{A reader adapted under the failure rule} \label{app:p6}

\begin{itemize}
\item \code{s110\_summary.py} compares the checksums carried by each result with the scripts in the folder where it
  runs, under their registered names \code{s110\_final.py} and \code{s110\_localise.py}. An adapted script therefore
  keeps its registered file name and replaces the original in the copy of this deposit that is used for the Sector 110
  runs.
\item The adapted photometric script is used unchanged for the primary and for every variant of that product. This is
  not an edit of a variant: variants are never edited individually.
\item The helper modules in \code{deps/} are never adapted. A pixel file that \code{deps/pixloc\_lib.py} cannot read is
  the case ``the target pixel file could not be read'' of Table~\ref{tab:failures}.
\item The sha256 of the adapted script, its diff against the deposited script, and the check of
  Appendix~\ref{app:procedure}, point 4, are posted before the Sector 110 run. A result whose \code{script\_sha256} is
  neither the checksum in \code{SHA256SUMS} nor a checksum posted in this way is invalid.
\item The summary compares only with the file on its own disk, so it accepts any script placed there; the check that
  counts is the public comparison of Appendix~\ref{app:procedure}, point 7.
\end{itemize}

\subsection{What ``robust'' means, and the line of variant outcomes} \label{app:p7}

\begin{itemize}
\item In the sentences, ``robust'' means exactly this: among the variants (a) to (e) that were run and reached
  $\Sens \geq 0.95$, none returned the opposite outcome (O3 against a primary O1; O1 or O2 against a primary O3). It
  does \emph{not} mean that all variants agree. A qualifying variant that returns O2 (a dip at the wrong depth or time)
  or O4 against a primary O1, or O4 against a primary O3, leaves ``robust'' in place. Variants below 0.95, variant (f),
  and variants that are not available or pending are not part of it.
\item Every posted summary is followed, for each signal, by one fixed line written by hand from the posted variant
  results and the summary file (the script does not write it). An example: ``Variant outcomes for signal A: a O1 (0.97);
  b O2, timing off (0.99); c O1 (0.96); d pending; e not available; f O1 (0.95; comparison only).'' The line lists every
  applicable variant in the order a to f. For each it gives the label O1, O2, O3 or O4 taken from the start of the
  outcome string of the script (for O2 followed by the failed conditions written by the script, for O4 by its reason)
  and its sensitivity to two decimals where its output file has one, else ``no sensitivity computed''; or ``not
  available''; or ``pending''. Variant (f) is listed in the same way, marked ``comparison only''. Wherever a sentence
  containing ``robust'' is quoted, this line is quoted with it.
\item ``Robust'' is never rewritten as ``confirmed by all processings'' or in any wider sense.
\end{itemize}

\section{How the results will be described} \label{app:wording}

The outcome labels are computed by \code{s110\_final.py} and \code{s110\_localise.py}. The sentences that describe them
are produced by \code{s110\_summary.py} from their output files, and only these sentences describe the outcomes.
Combined positions, variants, later reprocessings or products, and individual transit times and depths appear only in
tables, never in these sentences. ``$p$'' is always the empirical $p$ of Appendix~\ref{app:localisation}.
``Recovered'' never means ``confirmed planet'' (Section~\ref{sec:limits}). A combined statement about A and B, or about
the system, is not part of this test; it is exploratory and cannot be cited for or against either outcome. In a dry
run, ``Sector 110'' is replaced by ``the archival dry run (Sector NN)''.

In these tables X stands for the signal (A or B), and the other letters (n, R, k, N, U, r, u, C, s, P, x) are filled in
by the script: for example n is the number of counted transits, R and r are false-detection levels with their trial
counts k of N, U and u their bounds, s the sensitivity, P the empirical $p$, and N in the localisation part the number of
calibration injections. Each sentence about signal X is built from three parts: the outcome (Table~\ref{tab:sent1}), the robustness
qualifier for O1 and O3 (Table~\ref{tab:sent2}), and the localisation for O1 and O2 (Table~\ref{tab:sent3}).

\begin{longtable}{>{\raggedright\arraybackslash}p{0.08\textwidth}>{\raggedright\arraybackslash}p{0.86\textwidth}}
\caption{First part of the sentence, by outcome.\label{tab:sent1}}\\
\toprule
Outcome & First part of the sentence \\
\midrule
\endfirsthead
\toprule
Outcome & First part of the sentence \\
\midrule
\endhead
O1 & ``Signal X was recovered in Sector 110 at the pre-registered times and depth (n counted transits; measured
  false-detection level for this product about R \% overall (k of N trials; binomial 95 \% upper bound U \%, approximate
  because the trials share data) and r \% for exactly n counted transits (k' of N'; binomial bound u \%))'', with the
  numbers of Table~\ref{tab:fa} for the product of the run. If no trial had exactly $n$ transits: ``; no trials with
  exactly n counted transits''. For a TESS-SPOC primary: ``; measured on SPOC 2-min data, no trials exist for the
  TESS-SPOC product''. \\
O2 & ``A dip was found at the predicted times of signal X, but it is not the predicted signal (C), as pre-registered;
  this contradicts the archival model of signal X'', with C the failed conditions as written by the script: ``shallower
  than predicted'', ``deeper than predicted'', ``timing off''. \\
O3 & ``The predicted signal X is excluded by Sector 110 (pre-registered sensitivity s $>$= 0.95)'' \\
O4 & ``Sector 110 cannot decide on signal X (reason: R).'' with R one of the fixed reasons written by the script: fewer
  than 10 usable cadences; no registered QLP flux column; no counted transit; noise pool has fewer than 100 windows; d not
  above det and not below lo; sensitivity below 0.95; both of the last two. \\
\bottomrule
\end{longtable}

\begin{longtable}{>{\raggedright\arraybackslash}p{0.38\textwidth}>{\raggedright\arraybackslash}p{0.56\textwidth}}
\caption{Added part for the robustness qualifier (O1 and O3 only).\label{tab:sent2}}\\
\toprule
Case & Added part \\
\midrule
\endfirsthead
\toprule
Case & Added part \\
\midrule
\endhead
at least one variant qualifies, and none contradicts & ``; robust'' \\
no variant reaches $\Sens \geq 0.95$ & ``; robustness not testable'' \\
either of the two above while variant (d) is pending & the same, followed by `` (provisional until 2027-02-28)'' \\
O1, and a qualifying variant returned O3 & ``; a processing variant with adequate sensitivity excluded it (variant x);
  this is reported and does not change the outcome'' \\
O3, and a qualifying variant returned O1 or O2 & ``; a processing variant with adequate sensitivity detected a dip
  (variant x); this is reported and does not change the outcome'' \\
\bottomrule
\end{longtable}

\begin{longtable}{>{\raggedright\arraybackslash}p{0.22\textwidth}>{\raggedright\arraybackslash}p{0.72\textwidth}}
\caption{Added part for the localisation (O1 and O2 only).\label{tab:sent3}}\\
\toprule
Case & Added part \\
\midrule
\endfirsthead
\toprule
Case & Added part \\
\midrule
\endhead
A, L1 & ``; its pixel position is compatible with TIC 4206066 (empirical p = P from N on-target injections, not below
  0.05).'' \\
A, L2 & ``; its pixel position is in tension with TIC 4206066 (empirical p = P from N on-target injections, below 0.05
  and not below 0.0027).'' \\
A, L3 & ``; its pixel position is off target (empirical p = P from N on-target injections, below 0.0027): the signal
  most likely comes from another source, and a planet on TIC 4206066 is disfavoured.'' \\
A, not possible & ``; its position could not be measured (reason as written by the localisation script).'' \\
B, measured & ``; B's pixel offset is reported in the table without a label and is not interpreted in this test.'' \\
B, not possible & ``; B's localisation was not possible (reason as written by the localisation script).'' \\
A or B, fallback primary & as ``not possible'', with the reason ``fallback product: localisation not possible'' \\
\bottomrule
\end{longtable}

For O3 and O4 no localisation is run, and the sentence ends with the robustness part or the reason. P is printed with
four significant digits; the stated relation to the thresholds is computed from the unrounded value.

Two cases are worded by the code slightly differently from Tables~\ref{tab:sent1} and \ref{tab:sent3}; as stated in
Section~\ref{sec:protocol}, the code defines the sentence. The O4 reason ``both of the last two'' is written as ``d not
above det and not below lo; sensitivity below 0.95''. If the localisation distance is not finite, the label for A
is \code{NOT\_POSSIBLE}; the sentence for A then ends ``its position could not be measured (label NOT\_POSSIBLE).'',
and the sentence for B gives the clause ``B's pixel offset is reported in the table without a label and is not
interpreted in this test.''.

Appendix~\ref{app:public} adds three things: for the failure cases that leave no output file (Appendix~\ref{app:p5}) the
same sentence parts are written from the posted files with the fixed reasons listed there; a result declared invalid
and not repairable gives the reason ``pre-registered procedure not followed'' (Appendix~\ref{app:p1}); and every summary
is posted with the line of variant outcomes of Appendix~\ref{app:p7}.

\section{Exact commands} \label{app:commands}

All commands are run from the root of the deposit (Appendix~\ref{app:procedure}, point 2). Angle brackets mark the
names of the Sector 110 files.

\runin{Time-system check} (Appendix~\ref{app:p3}), before every run on a FITS file:
\begin{verbatim}
python -c "from astropy.io import fits; h = fits.getheader('FILE', 1); \
print({k: h.get(k, 'ABSENT') for k in ('TIMESYS', 'TIMEREF', 'BJDREFI', \
'BJDREFF', 'TIMEUNIT', 'TIMEDEL')})"
\end{verbatim}
Long commands are split over lines with a trailing backslash, so they can be pasted into a shell as they are.
Text in angle brackets is replaced by the actual file name or value.

\runin{Primary run:}
\begin{verbatim}
python s110_final.py --lc <tessNNNN-s0110-0000000004206066-NNNN-s_lc.fits> \
  --product spoc --out s110_primary.json
\end{verbatim}

\runin{Fallback 1} (only under Appendix~\ref{app:data}):
\begin{verbatim}
python s110_final.py --lc <QLP s0110 llc.fits> --product qlp --out s110_primary.json
\end{verbatim}
The flux column is \code{SYS\_RM\_FLUX}, or \code{DET\_FLUX} if the file has no \code{SYS\_RM\_FLUX}. If it has
neither, the script writes O4 for both signals with the reason ``no registered QLP flux column''.

\runin{Fallback 2:}
\begin{verbatim}
python s110_final.py --lc <TESS-SPOC s0110 lc.fits> --product spoc --out s110_primary.json
\end{verbatim}

\runin{Variants,} on the same product as the primary: add \code{--window 0.5} (a), \code{--window 1.0} (b),
\code{--drop-desat-hours 1} (c, SPOC primary only), \code{--flux SAP\_FLUX} (e), \code{--count-rule v2} (f); for (d)
use \code{--lc <QLP s0110 llc.fits> --product qlp} (SPOC primary only).

\runin{Localisation:}
\begin{verbatim}
python s110_localise.py --tpf <tessNNNN-s0110-0000000004206066-NNNN-s_tp.fits> \
  --photometry s110_primary.json --primary-sha256 <posted hash> \
  --prf-dir <MAST prf_fitsfiles/start_s0004> --out s110_loc_A.json
\end{verbatim}
and the same with \code{--signal B --out s110\_loc\_B.json}.

\runin{Sentences:}
\begin{verbatim}
python s110_summary.py --primary s110_primary.json --variant a=... \
  --loc-A s110_loc_A.json --loc-B s110_loc_B.json --d-status pending|final \
  [--d-final-reason crashed|not-public-on-deadline] --out s110_summary.json
\end{verbatim}
In the summary command, \code{--variant a=...} is given for every applicable variant, or
\code{--not-available KEY=FILE} instead for a variant that is not available. FILE is the attempt record of Appendix~\ref{app:p4}, which contains the unedited
error output of the attempted run, or, for (d), the saved MAST search result; its sha256 and text are copied into the
summary.

\code{--d-status pending} is used while no QLP Sector 110 release is public; the qualifier then stays provisional, and
this option is refused on or after 2027 February 28. \code{--d-status final} is used once (d) has been run
(\code{--variant d=...}), or with \code{--not-available d=FILE} and \code{--d-final-reason crashed} (a public release
was attempted and failed) or \code{not-public-on-deadline} (FILE is the saved MAST search of 2027 February 28 or later;
refused before that date). It is automatic for a QLP primary.

The summary script refuses to write sentences if a variant is missing; is not the registered variant of this primary;
carries another script, ephemeris, expected-depth or helper-module checksum than this package (audit 6); or is declared
not available without a saved non-empty error output. It also refuses if a required localisation file is missing, was
not computed from this primary result or by this package's \code{s110\_localise.py}, or is not from a Sector 110 pixel
file of TIC 4206066.

\section{Validation on archival data} \label{app:validation}

The frozen scripts were run on the archival sectors (folder \code{dryrun/}). These dry runs check that the code does
what Sections~\ref{sec:times} and \ref{sec:test} and Appendices~\ref{app:processing} to \ref{app:localisation}
say. They are not evidence about Sector 110 and not an
independent test, because Sectors 32 and 98 contributed to the ephemeris and to the expected depths. Archival epochs
are outside the frozen Sector 110 tables, so their $\sigma_{\rm pred}$ comes from the ephemeris covariance at those
epochs (about 3.0 to 3.5\,min instead of 4.1). The prediction part of the archival timing tolerances is therefore slightly smaller than in
Sector 110. Table~\ref{tab:dryruns} lists the results; depths and thresholds are in ppm.

\begin{longtable}{>{\raggedright\arraybackslash}p{0.25\textwidth}>{\raggedright\arraybackslash}p{0.12\textwidth}>{\raggedright\arraybackslash}p{0.33\textwidth}>{\raggedright\arraybackslash}p{0.20\textwidth}}
\caption{Dry runs of the frozen code on archival data.\label{tab:dryruns}}\\
\toprule
Data & Settings & A & B \\
\midrule
\endfirsthead
\toprule
Data & Settings & A & B \\
\midrule
\endhead
\multicolumn{4}{l}{\textit{Photometric test}} \\
S32 SPOC 2-min & primary & O1; $n = 6$ (epochs $-573$, $-566$ not tested, see Section~\ref{sec:depths}); $d = 437$;
  $\ddet = 287$ (280 to 302); band 114 to 766; sensitivity 0.92; $\Delta t = +2.5$\,min, tolerance 14.7; 0 of 58 fake
  phases above $\ddet$ & O4; no counted transit \\
S32, ephemeris shifted by half a period, real transits cut & null test & O4; $n = 7$; $d = -106$; $\ddet = 278$;
  sensitivity 0.92 (below 0.95, so not O3) & O4; $d = 194$; $\ddet = 369$; sensitivity 0.83 \\
S98 QLP & fallback 1 & O1; $n = 8$ (epoch $-1$ not tested: no baseline); $d = 450$; $\ddet = 150$; band 223 to 659;
  sensitivity 1.00; $\Delta t = +1.0$, tolerance 11.7 & O1; $n = 2$; $d = 995$; $\ddet = 306$; sensitivity 0.949 \\
S98 QLP, v2 counting rule (no baseline rule, no edge zones) & variant (f) & O4; $n = 9$; $d = 509$; $\ddet = 652$;
  sensitivity 0.05 & O4; $\ddet = 2428$; sensitivity 0.01 \\
S32 window 0.5\,d / 1.0\,d / desaturation $\pm 1$\,h & variants (a) to (c) & O1 / O1 / O1 (sensitivity 0.93 / 0.88 /
  0.94) & O4 (no counted transit) \\
S32 SAP flux & variant (e) & O1; sensitivity 0.96 & O4 (no counted transit) \\
S32, v2 counting rule & variant (f) & O1; $n = 8$ (the two epochs next to B counted); $d = 426$; $\ddet = 280$ & O4;
  $n = 2$; sensitivity 0.10 \\
\midrule
\multicolumn{4}{l}{\textit{Pixel localisation}} \\
S32 pixels, full localisation pipeline (equal stacks: 6 fake transits in all 3600 trials, 31 phase slots in the
  300-trial clouds and 32 in the calibration cloud, no failed trial) & Appendix~\ref{app:localisation} & offset $(+5.46,
  +0.38)''$, cloud standard deviation $(3.2, 4.1)''$, $\delta = 2.18$, empirical $p = 0.156$ (468 of 3000 calibration
  trials at least as far; the Gaussian tail of the earlier rule gave 0.093) (dry run, no label) & ``not possible'' (B
  not detected) \\
S32 pixels, coverage of the empirical $p$: the same run with calibration seed 5005, each of its 3000 on-target trials
  scored against the registered 3000, and the reverse & \code{build\slashbreak localisation\_coverage.py},
  \code{calib\slashbreak localisation\_coverage.json} & below 0.05: 4.7\,\% and 5.4\,\% (nominal 5\,\%); below 0.0027: 0.37\,\% and
  0.23\,\% (nominal 0.27\,\%); the Gaussian tail on the same trials: 13\,\% and 5.4 to 6.0\,\% & \\
S32 pixels, re-derivation ephemeris, all transits & fidelity check & offset $(+4.79, +0.90)''$, Gaussian $p = 0.213$ as
  in item 11 (empirical $p = 0.211$), all 600 injection positions identical to item 11 & \\
S32 null-test photometry & localisation gate & ``not possible'' (A not detected) & \\
\midrule
\multicolumn{4}{l}{\textit{Refusals and fixed sentences}} \\
S32 files, run without \code{--dry-run} & refusal check & photometry and localisation both refused (not Sector 110 /
  TIC 4206066) & \\
S6 TESS-SPOC pixel file (camera 1, CCD 1), run without \code{--dry-run}, non-registered switches & refusal check &
  refused before the camera/CCD rule, nothing written (audit 5) & \\
Sentences from the dry-run outputs (with \code{--dry-run}: ``in the archival dry run (Sector NN)'') &
  \code{s110\_summary.py} & S32: ``recovered \ldots\ robust (provisional until 2027-02-28); its pixel position is
  compatible with TIC 4206066 (empirical p = 0.1563 from 3000 on-target injections, not below 0.05)''; null test:
  ``cannot decide (reason: sensitivity below 0.95)''; S98: ``recovered \ldots\ robustness not testable; its position
  could not be measured (fallback product \ldots)'' & S32: ``cannot decide (reason: no counted transit)''; S98:
  ``recovered \ldots\ 0.62\,\% for exactly 2 counted transits \ldots'' \\
Summary with a variant silently left out & refusal check & refused (``variants \ldots\ are neither given nor declared
  not available'') & \\
Summary with a variant declared not available without a saved error output & refusal check & refused (audit 5) & \\
Summary with a variant that carries another script checksum (fixture \code{dryrun/s32\_sap\_other\_version.json}: a copy
  of \code{s32\_sap.json} whose script checksum is set to zeros on purpose) & refusal check & refused (audit 6) & \\
Summary with variant (d) declared final as ``not public'' before 2027-02-28 & refusal check & refused (audit 6) & \\
\bottomrule
\end{longtable}

The dry runs show the following in advance.
\begin{itemize}
\item At the noise level of the Sector 32 2-min data, with 6 or 7 counted transits, the test detects A (sensitivity
  0.92) but cannot return O3 in the null test (0.92 is below 0.95): a non-detection would be O4, undecided.
\item In Sector 110 at most 8 transits of A can be counted in practice. Epoch 117 (BTJD 4371.121) needs data until about
  BTJD 4371.53 (2026-11-27 00:41 UTC at Earth) to clear its edge zone, and epoch 109 needs continuous data from about
  BTJD 4345.24 (2026-10-31 17:45 UTC at Earth). More counted transits raise the sensitivity; it is computed from the
  Sector 110 data, not assumed.
\item Signal B, with 2 or 3 transits, will most likely end as O4 unless Sector 110 is quieter than Sector 32.
\item Without the baseline rule and the edge zones (variant f), single windows next to data gaps dominate the noise in
  Sector 98, and the test cannot decide even for the real signal. This is why those rules exist.
\item On these pixels the Gaussian tail of the earlier rule would have called about 6\,\% of truly on-target signals
  ``off target''; the empirical $p$ keeps this near the nominal 0.27\,\%.
\item In the localisation dry run the combined line counts Sector 32 twice and has no meaning; it only checks the code.
  Because B is not detected in the Sector 32 dry run, the B localisation is not exercised by any dry run.
\end{itemize}

\section{Known limits, stated before the data: full list} \label{app:limits}

\begin{itemize}
\item The archival sectors were inspected, and one choice in the ephemeris (B modelled jointly with A,
  Section~\ref{sec:predictions}) was made after looking at them.
\item With the injection-calibrated errors, the archival transit times scatter less than expected ($\chi^2 = 8.3$ for
  20 degrees of freedom). The prediction uncertainty of A may therefore be somewhat too large, which makes the timing
  rule more lenient.
\item $D_{98}$ for A (450\,ppm), the archival centre from which $\dhi$ is built, is close to the primary expectation
  (442\,ppm). For A the consistency band is therefore set mostly by the noise and by the two expectation errors, not by
  an excess in the discovery sector.
\item The edge-zone rule was introduced during the adversarial audit of a first draft, after the archival dry runs had
  shown the artefacts at data gaps. Before the rule was fixed, two zone widths (6\,h and 9\,h; at that stage measured from the
  nearest data gap) were tried once on archival data to see whether the artefacts were the cause. Both gave the same
  result (Sector 32: $\ddet = 211$ and 214\,ppm against 236 without zones; Sector 98: $\ddet = 142$ and 141 against
  581; no real transit lost), so the trial could not favour one width. The adopted width is $W/2$ by the definition of
  the filter, and it scales with $W$ in the variants. The archival recovery of A and B (Appendix~\ref{app:validation}) is
  a validation made after the rule was fixed, not the reason for it.
\item The change of percentiles (Appendix~\ref{app:history}, row 3b) was also made after archival measurements, to bring
  the measured false-alarm level down. It makes detection stricter and widens the consistency band slightly. The percentiles were set once
  (99.9 and 0.1) and then measured; they were not tuned further. The measured levels (SPOC 2-min: A 0.47\,\%, B
  0.49\,\% overall, with the thresholds of the Sector 110 run since audit 6; Appendix~\ref{app:fa}) pool every trial set
  that was run, including audit 3's. An intermediate version of this document quoted a lower level for A from one set
  only and a grouping of $n$ chosen after seeing the counts; both were corrected after audit 4.
\item The equal-stack localisation scheme was checked on audit 4's synthetic Sector 110 pixel file (a 2-d gap and 8
  transits imposed). The on-target cloud has standard deviations of 3.7 to $3.9''$ (East) and 4.7 to $5.2''$ (North),
  26 to 27 phase slots and no failed trial, against 3.4 to $3.5''$ and 9.2 to $10.7''$ with the earlier
  phase-rejection scheme. This check was made outside this deposit on a synthetic file built from the Sector 32 pixels;
  its scripts and outputs are kept by me and are not part of this deposit, so it is given for information only.
\item Edge zones cost data: in a Sector 110 light curve with two or three segments, a transit has roughly a 6 to 10\,\%
  chance to fall in an edge zone and not be tested.
\item Neighbouring windows of the noise pool overlap (steps of $\Tfourteen/2$), so the pool has fewer independent
  windows than its size. This does not bias the noise distribution, but its far tail rests on few independent windows.
  This is why the false-detection levels are measured separately (Appendix~\ref{app:fa}).
\item The archival localisation inputs include Sector 98 injection trials that overlap signal B (113 of 300). Using only
  the clean trials would change the archival combined $p$ from 0.061 to 0.053. The headline values are used as
  published. They come from the item 11 recipe, whose fakes may sit next to data gaps; the Sector 110 clouds exclude
  such fakes and are narrower. For the same Sector 32 dry-run offset $(+5.46, +0.38)''$ the Gaussian $p$ was 0.29 with
  the first wrapper (unequal stacks, fakes anywhere), 0.23 with equal stacks but no edge-zone rule, and 0.093 with the
  final scheme; the empirical $p$ of the final rule is 0.156. The combined value mixes the two kinds of cloud and uses
  Gaussian tails, and it is descriptive only.
\item The localisation $p$ was changed after audit 5 from the Gaussian tail $\exp(-\delta^2/2)$ to the empirical $p$,
  because archival on-target injections showed that the Gaussian tail is too light (Sector 32: 12 of 300 below
  0.0027). The thresholds 0.05 and 0.0027 were not changed, and the method was fixed before its value for the Sector 32
  dry run was computed. For that offset it raises $p$ from 0.093 to 0.156: such offsets are less unusual among
  on-target injections than a Gaussian suggests. For moderate offsets the final rule is therefore more lenient than the
  Gaussian one, while a false ``off target'' for an on-target signal now occurs near the nominal rate
  (Appendix~\ref{app:validation}). On audit 4's synthetic Sector 110 pixel file (8 transits, a 2-d gap; a check made
  outside this deposit, kept by me, for information only) the same coverage check gives 4.5 and 5.6\,\% below 0.05 and 0.07 and 0.50\,\% below
  0.0027 (Gaussian tail: 7.4 to 7.6\,\% and 1.7 to 1.9\,\%). All injections of one sector share pixels and overlapping
  windows, so the empirical $p$ assumes that a real on-target signal behaves like the injections at random phases. Its
  resolution is $1/(N_{\rm cal}+1)$, or 1/3001 when all 3000 calibration trials succeed. Near 0.0027 only about 8 calibration trials lie further out, so its Monte Carlo error there is
  about a third of its value.
\item Box depths are lower than fitted transit depths (limb darkening, finite exposure). The expected and measured
  values here are both box depths, so they can be compared with each other, but not with fitted depths quoted
  elsewhere.
\end{itemize}


\section{Changes from the internal draft v2, and why} \label{app:history}

These changes were made using archival data before Sector 110. The table retains the reasons for choices
made after examining those data; it is part of the scientific provenance. Editorial chronology and package
checks are recorded in \code{CHANGELOG.md}. Row numbers are used elsewhere in the protocol and package.

The post hoc choice in the ephemeris of A described in Section~\ref{sec:ephem} (B modelled jointly instead of masked)
is decision D13 of the final re-derivation (2026 October 1), labelled post hoc there.

\begingroup\footnotesize
\begin{longtable}{>{\raggedright\arraybackslash}p{0.04\textwidth}>{\raggedright\arraybackslash}p{0.22\textwidth}>{\raggedright\arraybackslash}p{0.33\textwidth}>{\raggedright\arraybackslash}p{0.33\textwidth}}
\caption{Changes from draft v2.\label{tab:history}}\\
\toprule
\# & Draft v2 & Problem & Final rule \\
\midrule
\endfirsthead
\toprule
\# & Draft v2 & Problem & Final rule \\
\midrule
\endhead
1 & ephemeris of A from the original joint fit (first transit BTJD 4345.66039, BJD\_TDB calendar 2026-11-01 03:49:49) & the final re-derivation refitted all 22 archival transit times with injection-calibrated errors & new ephemeris (Section~\ref{sec:predictions}); all nine predicted times of A move about 3.0 min earlier (2.96 to 2.97 min); B is unchanged \\
2 & timing tolerance from sigma\_pred = 1.5 x 2.0 min for A & the calibrated prediction error is 4.1 min & sigma\_pred of A from the full covariance of the new ephemeris (4.10 to 4.15 min per epoch), no extra factor; B keeps 1.5 x its joint-fit sigma (4.06 to 4.11 min) because its errors are not injection-calibrated \\
3 & noise windows cut from a trend fitted with only the real predictions masked & the real windows are masked during detrending (the trend under them is interpolated), the noise windows were not, so the noise looked quieter than the tested statistic; on archival data the 0.3 \% level was really about 1 \% (S32) to 6 \% (S98). v1 and v2 had the same flaw & the noise pool is built from fake ephemerides, each detrended with its fake windows masked exactly like the real ones (Appendix~\ref{app:noise}) \\
3b & det and lo at the 99.7th and 0.3rd percentiles of the bootstrapped noise & with the corrected noise model the realised false-alarm level, measured on archival fake ephemerides, was still 0.8 \% (A) and 1.2 \% (B) (audit 2, with the earlier segment definition): the bootstrap of a finite pool of overlapping windows underestimates the tails & det at the 99.9th, lo at the 0.1st, hi at the 99.9th percentile; realised levels measured per product and per number of transits in Appendix~\ref{app:fa}. This makes detection stricter and widens the consistency band slightly \\
4 & expected depths 437 $\pm$ 30 (A) and 630 $\pm$ 59 ppm (B) from v1, where all in-transit points of a sector were pooled & the test statistic is the mean of single-transit depths, a different estimator, and the v1 errors were white-noise errors & expected depths recomputed with exactly the test statistic, counting rule, detrending and noise model on the archival sectors (Section~\ref{sec:depths}) \\
5 & a transit counted if 80 \% of its window has data & a transit on a sector-start ramp with no baseline after egress was counted; in the v2 Sector 98 dry run it raised the statistic from 445 to 503 ppm & in addition, at least 2 cadences on each side between 0.75 and 3 T14 from mid-transit (baseline rule ``T3'' of the re-derivation) \\
5b & no rule for data gaps & with the corrected noise model the largest noise windows all sat within 0.1 d of a data gap, where the centred detrending filter has data on one side only and its trend drifts (single windows up to +4000 ppm in Sector 98) & edge zones: for every detrending the series actually given to the detrender (with that run's transit masks removed) is split into the same continuous segments the detrender fits (gaps $>$ 0.3 d), and the first and last W/2 of each segment are not used (W = detrending window, 0.75 d for the primary, so 9 h). A transit in an edge zone, or whose mask joins a data gap into a break, is ``not tested'', never a miss. Same rule for real and noise windows, both signals, primary, fallbacks and variants (with each variant's own W). A first version took the segments from the unmasked light curve; audit 2 showed that a short gap next to a transit mask then escaped the rule \\
6 & sigma\_meas = bootstrap over the counted transits & with one transit the bootstrap returns 0 & sigma\_meas from injections of the measured depth into noise-pool windows, defined for any number of transits \\
7 & ``individual transit times are reported''; ``input sha256 recorded'' & the script did neither & the script reports the shift of each counted transit (not part of the outcome) and records input, code and dependency hashes and package versions \\
8 & localisation ``identical to the Sector 32 analysis'', code outside the package, events chosen by the pixel code itself, archival combination undefined & the method was not frozen & the code is in the package (Appendix~\ref{app:localisation}), the same code as the final re-derivation (comments and default paths made neutral on 2026-10-05); transits come from the primary photometric run; the wrapper refuses non-registered settings; the combination is fixed \\
9 & product rules without cut-off dates; crashes and code edits open-ended; robustness rule one-sided & each leaves a choice after the data & Appendices~\ref{app:data}, \ref{app:variants} and \ref{app:procedure} \\
9b & (audit 3) fake stacks of the localisation had fewer transits than the real stack & a smaller fake stack widens the error cloud and pushes the label toward ``compatible'' & every fake stack now has exactly as many transits as the real one (Appendix~\ref{app:clouds}) \\
9c & (audit 3) the fixed wording did not cover every outcome; the false-detection level of A was an average over products & descriptions and numbers left open after the data & \code{s110\_summary.py} produces the only sentences allowed (Appendix~\ref{app:wording}); realised levels per product and per n (Appendix~\ref{app:fa}) \\
9d & (audit 4) the equal-stack rule rejected phases until enough fakes fitted, which, with all transits counted and a data gap, kept few phases and forced noisy fakes next to gaps into every stack; the false-detection level of A came from a re-draw that left out audit 3's identical trial set & the error cloud inflated (mainly North, toward ``compatible''); the quoted level was too low & fakes now obey the photometric edge-zone rule and missing fakes are taken from further phases (no phase rejected), with a floor of 15 phase slots per cloud; every trial set that was run is pooled, levels quoted overall and for the exact n, with upper bounds \\
9e & (audit 5) localisation label from the Gaussian tail p = exp(-d\textasciicircum{}2/2); the camera/CCD rule ran before the origin checks; a variant could be declared ``not available'' with any text & on the archival Sector 32 pixels 12 of 300 on-target injections had exp(-d\textasciicircum{}2/2) $<$ 0.0027 (Gaussian expectation about 1), so the tail was not calibrated; an archival file from another sector on another CCD wrote an official ``not possible''; an unwanted variant could be dropped from the robustness rule & the label uses an empirical p from 3000 further on-target injections, thresholds unchanged (Appendix~\ref{app:localisation}); every origin check comes before the camera/CCD rule, and \code{s110\_summary.py} checks the sector and star of each localisation file; ``not available'' needs the saved error output, whose hash and text enter the summary \\
9f & (audit 6) a saved ``QLP not yet public'' record made the robustness qualifier final; the summary did not compare the script, ephemeris, expected-depth and helper-module hashes of variants and localisation files with the package; four localisation dry-run files carried the hash of an earlier text of the localisation script; the false-alarm trials used one seed where the Sector 110 run uses the median of 20; the fixed sentence wrote ``at most'' for a binomial bound of dependent trials & each left a way to describe Sector 110 differently from the rules, or quoted a level not measured with the registered procedure (on archival data one phase changed class with the median of 20 seeds) & ``not public'' keeps the qualifier provisional; only a run (d), a failed attempt on a public release, or a MAST search on or after 2027-02-28 makes it final (Appendix~\ref{app:robust}, checked by the script); every result must carry this package's hashes; the four files were re-run with the final script, all numbers identical; all trial sets were re-run with the thresholds of the Sector 110 run (14 of 17 835 classifications changed; Appendix~\ref{app:fa}); the bound is called an approximation (Appendices~\ref{app:fa} and \ref{app:wording}) \\
9g & (final independent check, 2026-10-03) the column of predicted times headed ``UTC'' gave the barycentric time as a calendar date; the text did not tie a fallback primary or a variant declared ``not available'' to saved evidence; the O4 row listed P(detect | real) $<$ 0.95 without ``d $\leq$ det''; the inputs of the ephemeris, a licence text and a complete reproduction recipe were missing & a reader could take the barycentric dates for arrival times (they are 5.6 to 7.2 min late), or read a dip above det with low sensitivity as O4 & the table gives the BJD\_TDB calendar date and the arrival time at the Earth in UTC (Section~\ref{sec:times}); Appendices~\ref{app:data}, \ref{app:variants} and \ref{app:procedure} tie fallbacks and ``not available'' to saved evidence; the O4 row states what the code does; \code{calib/transit\_times\_A.csv}, \code{calib/ephemeris\_inputs.json}, \code{LICENSE} and the reproduction recipe in \code{README.md} were added. No code, number or threshold changed \\
9h & (adversarial check of the deposit, 2026-10-03) the summary accepts any non-empty file as the error output of a variant declared not available; a failure that ends without a JSON (twice failed photometry, missing or unreadable pixel file) has no path to a fixed sentence; a reading adaptation under Appendix~\ref{app:procedure}, point 4, did not say under which file name it runs, so the summary could refuse it; ``robust'' also covers a sensitive variant returning O2; the code does not check that a fallback product was allowed when it was chosen, nor the time system of an input file & each left a way to use or describe a Sector 110 result differently from the rules without the code noticing & Appendix~\ref{app:public}: a result obtained in violation of a pre-registered rule is invalid; written attempt records with fixed content for every run and every ``not available''; fixed sentences, written from the posted files, for the failure cases without a JSON; an adapted reader keeps the registered file name; the exact meaning of ``robust'' and a fixed line listing every variant's outcome; a saved MAST product search before any fallback; registered time-system header values checked and posted before every run. The code was not changed and does not enforce these rules; the public record does. No code, number or threshold changed \\*
\bottomrule
\end{longtable}
\endgroup

\section{Files of the deposit} \label{app:files}

\begin{itemize}
\item \code{PREREG\_S110\_FINAL.pdf} (this document) and its LaTeX source in \code{protocol\_source/} (main
  file \code{PREREG\_S110\_FINAL.tex}): the authoritative text of the pre-registration (version 1.1). It replaces the
  earlier Markdown text \code{PREREG\_S110\_FINAL.md} of 2026 October 5, which is not part of this deposit; the
  changes are listed in \code{CHANGELOG.md}.
\item \code{PREREG\_S110\_SHORT.pdf}: a separate reader-oriented overview; this full protocol governs (source \code{protocol\_source/PREREG\_S110\_SHORT.tex}).
\item \code{CHANGELOG.md}, \code{README.md} (overview, reproduction recipe, acknowledgements), \code{FILE\_GUIDE.md} (the
  purpose of every file), \code{LICENSE}, \code{CITATION.cff}, \code{requirements.txt} (the binding pin),
  \code{requirements\_lock.txt} (the full environment with checksums), \code{SHA256SUMS}.
\item \code{calib/transit\_times\_A.csv}, \code{calib/ephemeris\_inputs.json} and \code{calib/ephemeris\_sources/}: the
  inputs of \code{ephemeris\_final.json} and their sources (added 2026 October 3). Every number in them is copied
  unchanged from my archival analyses. On 2026 October 5 only the free-text field \code{B.ephemeris\_note} of
  \code{calib/ephemeris\_inputs.json} was reworded; no number or other field changed.
\item \code{build/fetch\_archival\_inputs.sh}, \code{build/reproduce.sh}, \code{build/compare\_outputs.py}: they download
  the public archival inputs, rebuild every output in a copy of the deposit and compare it with this deposit (added
  2026 October 3; they do not change any analysis).
\item \code{s110\_final.py}, \code{s110\_localise.py}, \code{s110\_summary.py}, \code{ephemeris\_final.json};
  \code{deps/lcio.py} and \code{deps/pixloc\_lib.py} (frozen copies).
\item \code{calib/expected\_depths.json}, \code{calib/localisation\_archival.json};
  \code{calib/localisation\_coverage.json} and \code{build/localisation\_coverage.py} (coverage of the empirical
  localisation $p$, Appendix~\ref{app:validation}); \code{calib/false\_alarm\_*.json} (trials),
  \code{calib/false\_alarm\_levels.json} (levels), \code{calib/false\_alarm\_summary.txt},
  \code{build/false\_alarm\_check.py} and \code{build/false\_alarm\_levels.py} (measured false-alarm levels,
  Appendix~\ref{app:fa}).
\item \code{dryrun/} (outputs on archival data); \code{build/make\_ephemeris.py} (how the ephemeris file was built) and
  \code{build/run\_all.sh} (rebuilds \code{calib/expected\_depths.json} and every dry run).
\item The internal drafts v1 (2026 September 27) and v2 (2026 September 30) are not part of this deposit; they are kept
  unchanged by me (Section~\ref{sec:protocol}).
\item \code{PRF\_SHA256SUMS}: the sha256 checksums of the TESS PRF files used (MAST
  \code{prf\_fitsfiles/start\_s0004/cam1\_ccd2/}). They are public MAST files and are not copied here; the localisation
  output records their checksums on every run.
\end{itemize}

\section{Section numbers used in the code} \label{app:codemap}

The three \code{s110\_*.py} scripts were not changed for this version. Their comments, help texts and one error message refer to sections of the
earlier Markdown text \code{PREREG\_S110\_FINAL.md} (for example ``section 9b'' in \code{s110\_summary.py}, or the
message ``the status of variant (d) is final (section 3)''). Table~\ref{tab:codemap} gives the place of each of those
sections in the present document.

\begin{longtable}{>{\raggedright\arraybackslash}p{0.17\textwidth}>{\raggedright\arraybackslash}p{0.33\textwidth}>{\raggedright\arraybackslash}p{0.42\textwidth}}
\caption{Sections of the earlier text and their place in this document.\label{tab:codemap}}\\
\toprule
Earlier section & Subject & Place here \\
\midrule
\endfirsthead
\toprule
Earlier section & Subject & Place here \\
\midrule
\endhead
opening & what the pre-registration tests & Sections~\ref{sec:question} and \ref{sec:signals} \\
0 & what the document is & Sections~\ref{sec:question} and \ref{sec:protocol} \\
1 & changes from draft v2 & Appendix~\ref{app:history} \\
2 & data and dates & Appendix~\ref{app:data} \\
3 & ephemeris, predicted times, processing, statistic, variants, robustness & Sections~\ref{sec:ephem},
  \ref{sec:times} and \ref{sec:stat};
  Appendices~\ref{app:processing} and \ref{app:variants} \\
4 & decision rule, expected depths, false-detection levels & Sections~\ref{sec:depths} and \ref{sec:outcomes};
  Appendices~\ref{app:decision} and \ref{app:depthnotes} \\
5 & timing rule & Section~\ref{sec:timing} \\
6 & localisation rule & Section~\ref{sec:localisation}; Appendix~\ref{app:localisation} \\
7 & not tested & Section~\ref{sec:limits} \\
8 & dry runs on archival data & Appendix~\ref{app:validation} \\
9 & commitments and exact procedure & Appendices~\ref{app:procedure} and \ref{app:commands} \\
9b & fixed wording of the results & Appendix~\ref{app:wording} \\
9c, 9c.1 to 9c.7 & rules the code does not enforce & Appendix~\ref{app:public}, Appendices~\ref{app:p1} to \ref{app:p7} \\
10 & known limits & Section~\ref{sec:limits}; Appendix~\ref{app:limits} \\
10b & what Sector 110 cannot show & Section~\ref{sec:limits} \\
11 & files & Appendix~\ref{app:files} \\
\bottomrule
\end{longtable}


\bibliography{references}{}
\bibliographystyle{aasjournalv7}

\end{document}
